%% admissibility_principle.tex
%% The Admissibility Principle: Computation, Physics, and Analogy as One Operation
%% Flyxion (standardgalactic) — Independent Researcher
%% LuaLaTeX recommended; pdfLaTeX compatible with minor adjustments

\documentclass[12pt, a4paper]{article}

% -------------------------------------------------------
% Packages
% -------------------------------------------------------
\usepackage{fontenc}
\usepackage[utf8]{inputenc}
\usepackage{microtype}
\usepackage{geometry}
\geometry{a4paper, top=2.8cm, bottom=2.8cm, left=3cm, right=3cm}

\usepackage{amsmath, amssymb, amsthm}
\usepackage{mathtools}
\usepackage{amscd}          % commutative diagrams
\usepackage{stmaryrd}       % for semantic brackets \llbracket \rrbracket

\usepackage{booktabs}       % professional tables
\usepackage{array}
\usepackage{longtable}

\usepackage{hyperref}
\hypersetup{
  colorlinks = true,
  linkcolor  = black,
  citecolor  = black,
  urlcolor   = black,
  pdftitle   = {The Admissibility Principle},
  pdfauthor  = {Flyxion},
}

\usepackage{setspace}
\setstretch{1.15}

\usepackage{titlesec}
\titleformat{\section}{\large\bfseries}{\thesection.}{0.6em}{}
\titleformat{\subsection}{\normalsize\bfseries}{\thesubsection.}{0.5em}{}
\titleformat{\subsubsection}{\normalsize\itshape}{\thesubsubsection.}{0.5em}{}

\usepackage{enumitem}
\setlist[itemize]{leftmargin=1.5em, itemsep=0.2em}

\usepackage{tcolorbox}
\tcbuselibrary{skins, breakable}

% Epistemic-status boxes
\newtcolorbox{established}{
  colback=gray!8, colframe=gray!50, fonttitle=\small\bfseries,
  title=Epistemic Status: Established (peer-reviewed results),
  breakable, left=4pt, right=4pt, top=4pt, bottom=4pt
}
\newtcolorbox{prototype}{
  colback=blue!5, colframe=blue!35, fonttitle=\small\bfseries,
  title=Epistemic Status: Engineering Prototype / Proposed Architecture,
  breakable, left=4pt, right=4pt, top=4pt, bottom=4pt
}
\newtcolorbox{speculative}{
  colback=orange!6, colframe=orange!45, fonttitle=\small\bfseries,
  title=Epistemic Status: Speculative Theoretical Proposal,
  breakable, left=4pt, right=4pt, top=4pt, bottom=4pt
}

% Theorem environments
\theoremstyle{plain}
\newtheorem{theorem}{Theorem}[section]
\newtheorem{proposition}[theorem]{Proposition}
\newtheorem{lemma}[theorem]{Lemma}
\newtheorem{corollary}[theorem]{Corollary}

\theoremstyle{definition}
\newtheorem{definition}[theorem]{Definition}
\newtheorem{example}[theorem]{Example}
\newtheorem{schema}[theorem]{Schema}

\theoremstyle{remark}
\newtheorem{remark}[theorem]{Remark}
\newtheorem{note}[theorem]{Note}

% Convenience macros
\newcommand{\betal}{\to_{\beta}}
\newcommand{\betastar}{\to_{\beta}^{*}}
\newcommand{\Lam}[2]{\lambda #1.\, #2}
\newcommand{\App}[2]{(#1\, #2)}
\newcommand{\subst}[3]{#1[#2 \mathrel{:=} #3]}
\newcommand{\Rcal}{\mathcal{R}}
\newcommand{\Bcal}{\mathcal{B}}
\newcommand{\Acal}{\mathcal{A}}
\newcommand{\Dcal}{\mathcal{D}}
\newcommand{\Phi}{\Phi}
\newcommand{\nats}{\;\text{nats}}

% -------------------------------------------------------
\begin{document}
% -------------------------------------------------------

% Title page
\begin{titlepage}
  \centering
  {\LARGE\bfseries The Admissibility Principle}\\[0.5em]
  {\large\itshape Computation, Physics, and Analogy as One Operation}\\[2.0cm]
  {\large Flyxion}\\[0.4em]
  {\normalsize Independent Researcher}\\[0.4em]
  {\normalsize May 2026}\\[1.0cm]
  \begin{abstract}
    \noindent
    We propose that a single abstract operation underlies
    $\beta$-reduction in the $\lambda$-calculus,
    membrane collapse in the Spherepop computational geometry,
    analogy-driven chunking in Hofstadter's theory of cognition,
    cellular differentiation as reconstructed through multimodal lineage
    tracing (Wang, He, \& Hu 2026),
    entropy-constrained field relaxation in the Relativistic Scalar-Vector
    Plenum (RSVP) framework,
    quantum interference-mediated dwell-time anomalies,
    phase-matched energy redistribution in ultrafast nonlinear optics,
    maximum-dissipation selection in turbulent fluid dynamics,
    recursive epistemic admissibility in strategic games,
    and temporal-domain feature extraction in the Marine signal
    architecture.
    We call this operation \emph{admissibility-constrained collapse}: a
    bounded region accepts a compatible perturbation, internally
    reorganizes its constraint structure, and collapses into a new
    admissible successor state while preserving certain invariants.
    The paper formalizes a generalized reduction operator $\Rcal$ of which
    each domain provides a specialization, and identifies a family of
    local-compatibility reconciliation structures --- exemplified by
    Church--Rosser confluence, sheaf gluing, and iterated admissibility
    elimination --- that recur across domains as instances of a shared
    convergence grammar rather than a single proved universal theorem.
    Spherepop is situated as a proposed geometric operational semantics
    mediating between the symbolic ($\lambda$-calculus), cognitive
    (Hofstadter), and thermodynamic (RSVP) levels, and its current status
    as a research program awaiting axiomatization is acknowledged
    explicitly.
    The biological case of cellular differentiation is identified as the
    strongest empirical pillar: lineage reconstruction requires historical
    admissibility propagation as a computational necessity, not as a
    theoretical overlay.
    Epistemic status is stratified throughout into five explicit layers:
    established peer-reviewed results, engineering prototypes, formal
    abstraction, philosophical synthesis, and open research program.
  \end{abstract}
\end{titlepage}

\tableofcontents
\newpage

% ================================================================
\section{Introduction and Motivation}
% ================================================================

\subsection{The central claim}

What does a photon do when it briefly becomes an atomic excitation?
What does a $\lambda$-term do when it swallows its argument?
What does a mind do when a new perception activates a stored conceptual chunk?
What does a nonlinear crystal do when a chirped pulse traverses it?
What does a pluripotent stem cell do when it commits to a developmental lineage?

Superficially these are questions from different sciences with different
vocabularies.
We argue they share a single formal skeleton: a bounded region equipped with
a constraint structure accepts an incoming perturbation, tests its local
compatibility, and collapses into a new configuration while preserving an
invariant.
We call this the \emph{admissibility principle}.

The principle has at least four realizations that can be made formally
precise and at least three more that we treat as motivated engineering
architectures or speculative proposals. These realizations are not
physically identical, and the paper is careful not to conflate different
evidential strata.
The unity we claim is structural, not material: the same abstract grammar
$\Rcal(\Bcal_i, \Delta_i) \to \Bcal_{i+1}$ governs each case at an
appropriate level of abstraction.

\subsection{Why this matters}

Standard cross-disciplinary syntheses suffer one of two failure modes.
The first is \emph{premature literalism}: claiming that disparate systems
are physically the same or that metaphors constitute proofs.
The second is \emph{vacuous analogy}: the observation that ``everything
transforms'' has no predictive content.

The admissibility principle attempts to occupy the productive middle ground
by requiring, for each domain, an explicit identification of:
\begin{itemize}
  \item the admissibility region $\Bcal_i$ and its constraint structure;
  \item the perturbation $\Delta_i$ and its compatibility condition;
  \item the collapse operator $\Rcal$ and what it preserves;
  \item the invariant or normal form $\Bcal_{i+1}$.
\end{itemize}
Where these four elements can be made precise, the analogy is a theorem
or a well-defined conjecture.
Where they cannot, the analogy remains heuristic and is labeled accordingly.

\subsection{What this framework does not claim}

It is important to state explicitly what the admissibility principle
is \emph{not}.
\begin{itemize}
  \item It does not claim that cognition is literally $\lambda$-calculus,
        or that turbulence is secretly semantics, or that RSVP parameters
        are established physical constants.
  \item It does not claim that all admissibility systems are
        mathematically identical or that the generalized reduction operator
        $\Rcal$ is a proven theorem rather than a research schema.
  \item It does not claim that every constrained process instantiates
        $\Rcal$.
        A thermostat constrains temperature but does not qualify: it
        filters outputs without recursively restructuring the state space.
        A lookup table maps inputs to outputs without propagating a
        transportable invariant.
        Mere constraint does not entail admissibility-collapse.
  \item It does not claim formal equivalence between domains that are
        only structurally analogous.
        The paper distinguishes throughout: formal equivalence
        (provable), structural analogy (argued), admissibility-preserving
        transport (conjectured), and metaphorical correspondence
        (labeled as such).
\end{itemize}

What the framework \emph{does} claim is that many systems exhibit
recurring structural motifs involving constrained trajectory reduction,
invariant preservation, and local compatibility propagation, and that
these motifs are sufficiently precise to motivate a shared research
grammar.
The analogy to category theory is apt: category theory did not claim
all mathematics was literally one object, but supplied a transport
language for structural relationships across domains.
The admissibility principle aims to play the same role for
constraint-preserving collapse.

\subsection{Organization}

The paper is organized into four parts plus a discussion, followed by
seven appendices.

\textbf{Part~I} (Sections~\ref{sec:lambda}--\ref{sec:analogy-reduction})
develops the formally precise core: $\lambda$-calculus and the
Curry--Howard correspondence, Spherepop as proposed geometric operational
semantics, Hofstadter's theory of analogy as the cognitive reading, and
the formulation of analogy itself as a reduction operation.

\textbf{Part~II} (Section~\ref{sec:lineage}) provides the biological
instantiation: cellular differentiation reconstructed through multimodal
lineage tracing, drawn from Wang, He, and Hu (2026) in
\textit{Nature Reviews Genetics}.
This is the domain in which admissibility is a computational necessity
rather than an interpretive overlay, and it serves as the strongest
empirical pillar of the monograph.

\textbf{Part~III} (Sections~\ref{sec:quantum}--\ref{sec:optics})
treats established physics: quantum weak-value dwell times and DC-OPA
ultrafast optics, each mapped onto the admissibility schema.

\textbf{Part~IV} (Sections~\ref{sec:rsvp}--\ref{sec:phoenix})
develops the speculative and prototypical layers: the RSVP thermodynamic
framework, the Marine/Phoenix 8a signal architecture, and the MEM$|$8
semantic memory proposal, all explicitly marked as conjectural or
engineering-stage.

\textbf{Part~V} (Sections~\ref{sec:unification}--\ref{sec:rsvp-synthesis})
synthesizes all layers into the generalized reduction operator, develops
the game-theoretic, turbulence, and relativistic admissibility cases,
articulates the five-layer epistemic architecture, and integrates RSVP,
CLIO, TARTAN, and Yarncrawler.

Section~\ref{sec:conclusion} identifies open problems, distinguishes what
has been established from what remains to be done, and states the
framework's self-assessment as a research grammar rather than a unified
theory.
Section~\ref{sec:underconstraint} develops the underconstraint problem
following Yudkowsky (2008) and connects it to the sheaf-theoretic
obstruction framework.
The seven appendices develop the category-theoretic, fiber-bundle,
sheaf-cohomological, variational, homotopy-theoretic, non-vacuity, and
three-level coarse-graining formulations in mathematical detail.

% ================================================================
\part{The Formal Core}
% ================================================================

\setcounter{section}{1}

% ================================================================
\section{$\lambda$-Calculus, $\beta$-Reduction, and Curry--Howard}
\label{sec:lambda}
% ================================================================

\begin{established}
  The results in this section are standard and fully established.
  Primary sources: Church (1936), Curry and Feys (1958),
  Howard (1980), Girard (1986), and Sørensen--Urzyczyn (2006).
\end{established}

\subsection{Syntax and reduction}

The untyped $\lambda$-calculus is defined over the grammar
\[
  M, N \;\Coloneqq\; x \mid \Lam{x}{M} \mid \App{M}{N},
\]
where $x$ ranges over a countably infinite set of variables.
The canonical computation rule is \emph{$\beta$-reduction}:
\begin{equation}
  \App{(\Lam{x}{M})}{N} \;\betal\; \subst{M}{x}{N}.
  \label{eq:beta}
\end{equation}
The substitution $\subst{M}{x}{N}$ replaces every free occurrence of $x$
in $M$ with $N$, with the standard capture-avoidance proviso
(renaming bound variables, $\alpha$-conversion, as necessary).

\begin{definition}[Admissibility in $\lambda$-calculus]
  A redex $\App{(\Lam{x}{M})}{N}$ is \emph{admissible} if $N$ is
  free for $x$ in $M$ modulo $\alpha$-conversion.
  Every redex is admissible after appropriate renaming.
\end{definition}

The reduction relation $\betastar$ is the reflexive transitive closure of
$\betal$.
A term $M$ is in \emph{$\beta$-normal form} if it contains no redex.

\subsection{Church--Rosser confluence}

\begin{theorem}[Church--Rosser, 1936]
  \label{thm:cr}
  If $M \betastar M_1$ and $M \betastar M_2$, then there exists $M_3$
  such that $M_1 \betastar M_3$ and $M_2 \betastar M_3$.
\end{theorem}

Theorem~\ref{thm:cr} asserts that distinct reduction paths from a
common source can always be \emph{joined}: there exists a common
reduct downstream.
Normal forms, when they exist, are therefore unique up to
$\alpha$-equivalence.
This is the first instance of the master convergence principle that we
will trace across all domains.

\subsection{Curry--Howard: types as propositions}

The \emph{Curry--Howard correspondence} identifies:
\begin{align*}
  \text{propositions} &\;\leftrightarrow\; \text{types,}\\
  \text{proofs}       &\;\leftrightarrow\; \text{typed terms,}\\
  \text{proof normalization} &\;\leftrightarrow\; \text{$\beta$-reduction,}\\
  \text{provability}  &\;\leftrightarrow\; \text{type inhabitation.}
\end{align*}
Under the simply-typed $\lambda$-calculus, every type is a proposition
of minimal implicational logic and every well-typed term is a proof.
Normalization of proofs corresponds exactly to $\beta$-reduction of terms.

\begin{remark}
  The BHK (Brouwer--Heyting--Kolmogorov) interpretation provides the
  semantic backing: a proof of $A \Rightarrow B$ is a \emph{procedure}
  that transforms proofs of $A$ into proofs of $B$.
  This already anticipates the admissibility principle:
  a type/proposition is a constraint, a term/proof is the compatible
  perturbation, and normalization is the collapse into normal form.
\end{remark}

\subsection{Divergence and fixed points}

In the untyped setting, not all terms normalize.
The combinator
\[
  \Omega \;=\; \App{(\Lam{x}{\App{x}{x}})}{(\Lam{x}{\App{x}{x}})}
\]
diverges: $\Omega \betal \Omega \betal \cdots$.
The fixed-point combinator
\[
  Y \;=\; \Lam{f}{\App{(\Lam{x}{\App{f}{\App{x}{x}}})}{(\Lam{x}{\App{f}{\App{x}{x}}})}}
\]
satisfies $\App{F}{(\App{Y}{F})} =_\beta \App{Y}{F}$ for any $F$.
Divergence in our schema corresponds to an admissibility loop with no
stable attractor: the system perpetually reorganizes without collapsing
to a normal form.

% ================================================================
\section{Spherepop: Geometric Operational Semantics}
\label{sec:spherepop}
% ================================================================

\begin{speculative}
  Spherepop is a theoretical computational framework under development.
  The correspondence with $\lambda$-calculus stated here is a formal
  proposal, not an independently verified result.
  The claim is that the mapping in Table~\ref{tab:lambda-spherepop}
  constitutes a sound interpretation; proving this requires a complete
  formal specification of Spherepop syntax and reduction rules,
  which is ongoing work.
\end{speculative}

\subsection{Motivation}

Standard $\lambda$-calculus is syntactic: computation is symbol
manipulation governed by string-rewriting rules.
Substitution $\subst{M}{x}{N}$ is a meta-theoretic operation,
not an object-level entity.

Spherepop proposes to \emph{externalize} substitution as a spatial,
topological, and historical operation.
The guiding intuition: a $\lambda$-abstraction creates a
\emph{protected evaluation region}; application injects an argument
into that region; substitution propagates through it;
the region boundary dissolves once computation is resolved.
This sequence is not a metaphor — it is a proposed formal
operational semantics in which boundaries are first-class objects.

\subsection{Core vocabulary}

\begin{definition}[Bubble]
  A \emph{bubble} $\mathbf{b} = (I, \partial, H)$ consists of:
  \begin{itemize}
    \item an \emph{interior} $I$: a set of latent sub-expressions;
    \item a \emph{membrane} $\partial$: a constraint structure
          specifying what perturbations are locally admissible;
    \item a \emph{history} $H$: an ordered record of prior reductions
          that produced the current state.
  \end{itemize}
\end{definition}

\begin{definition}[Admissible injection]
  A bubble $\mathbf{b}'$ is \emph{admissible} for injection into
  $\mathbf{b}$ if $\mathbf{b}'$ is compatible with $\partial(\mathbf{b})$
  in the sense that all free attachment sites in $I(\mathbf{b})$
  can be bound by the structure of $\mathbf{b}'$ without capturing
  existing bindings in $H(\mathbf{b})$.
\end{definition}

\begin{definition}[Membrane collapse]
  Given an admissible injection of $\mathbf{b}'$ into $\mathbf{b}$,
  \emph{membrane collapse} is the operation:
  \[
    \mathbf{b} \oplus \mathbf{b}' \;\longrightarrow\; \mathbf{b}''
  \]
  where $\mathbf{b}''$ has:
  \begin{itemize}
    \item interior $I(\mathbf{b}'') = \subst{I(\mathbf{b})}
          {\text{open sites}}{\mathbf{b}'}$,
    \item history $H(\mathbf{b}'') = H(\mathbf{b}) \cdot \langle\mathbf{b}'\rangle$
          (history extended with the injection event),
    \item membrane $\partial(\mathbf{b}'')$ determined by the residual
          open sites of $I(\mathbf{b}'')$.
  \end{itemize}
\end{definition}

The key distinction from standard substitution: $H$ makes collapse
\emph{irreversible} and \emph{identity-preserving}.
Two bubbles with identical interiors but different histories are
distinct computational objects.

\subsection{Correspondence with $\lambda$-calculus}

\begin{proposition}[Spherepop--$\lambda$ correspondence]
  Under the interpretation of Table~\ref{tab:lambda-spherepop},
  membrane collapse is the geometric realization of $\beta$-reduction.
  $\alpha$-conversion corresponds to historical rebinding of attachment
  sites; confluence corresponds to path-independent collapse producing
  the same stable topology.
\end{proposition}

\begin{table}[h!]
\centering
\caption{Correspondence between $\lambda$-calculus and Spherepop.}
\label{tab:lambda-spherepop}
\begin{tabular}{ll}
\toprule
\textbf{$\lambda$-calculus} & \textbf{Spherepop} \\
\midrule
Variable          & Open attachment site \\
Abstraction       & Scoped bubble with membrane \\
Application       & Bubble injection \\
$\beta$-reduction & Membrane collapse \\
$\alpha$-conversion & Historical rebinding \\
$\beta$-normal form & Stable bubble configuration \\
Divergence ($\Omega$) & Infinite recursive bubbling \\
Confluence (Church--Rosser) & Path-independent topological collapse \\
\bottomrule
\end{tabular}
\end{table}

\subsection{Spherepop as a research program rather than a completed formalism}

The reviewer of any draft of this monograph will correctly observe that
Spherepop currently lacks a complete operational semantics: there are no
published reduction rules, no formal grammar with a proved decidable
word problem, no normalization theorem.
This observation is accurate and should be acknowledged directly.

Spherepop currently occupies the historical position that manifold
theory occupied before rigorous atlases and transition maps were
axiomatized, or that tensor notation occupied before Ricci calculus was
stabilized.
Long before these formalisms were complete, geometers and physicists
reasoned productively through intuition and local patching.
Spherepop is at that pre-axiomatic stage.

What the present work contributes is the identification of the minimal
ingredients that a completed Spherepop formalism would require:
\begin{enumerate}
  \item \emph{Scoped regions} with formally defined interiors and
        membranes;
  \item \emph{Historical identity persistence} distinguishing bubbles
        with identical interiors but different causal ancestry;
  \item \emph{Admissible injection}, defined by a compatibility
        predicate on membranes;
  \item \emph{Collapse propagation}, specifying how substitution
        propagates through the interior after injection;
  \item \emph{Topology-sensitive rebinding}, the analogue of
        $\alpha$-conversion in the spatial setting.
\end{enumerate}
These ingredients correspond exactly to the components of explicit
substitution calculi ($\lambda\sigma$, $\lambda\upsilon$), with the
addition of irreversible historical extension.
Formalizing Spherepop is therefore a well-posed research problem,
not an indefinitely open aspiration.

The present paper treats Spherepop as a proposed geometric semantics
awaiting axiomatization, not as a completed system.
Claims involving Spherepop are accordingly marked
\emph{Epistemic Status: Speculative Theoretical Proposal} throughout.

\subsection{Explicit substitution and visible dependency}

Traditional $\lambda$-calculus hides substitution as a meta-operation.
Explicit substitution calculi (de Bruijn indices, $\lambda\sigma$, $\lambda\upsilon$)
bring it into the object language.
Spherepop can be understood as a \emph{spatial explicit substitution calculus}:
every bubble carries a visible dependency graph describing what may enter,
what is bound, and which historical trajectories remain admissible
after collapse.
This makes Spherepop closer in spirit to
$\lambda\upsilon$ or $\lambda x$ calculi than to the original Church system,
but with the additional dimension of irreversible event history.

% ================================================================
\section{Hofstadter: Analogy as Admissibility-Driven Collapse}
\label{sec:hofstadter}
% ================================================================

\begin{established}
  Hofstadter's theoretical claims about analogy-making are from
  published work (\textit{Gödel, Escher, Bach}, 1979;
  \textit{Fluid Concepts and Creative Analogies}, 1995;
  \textit{Surfaces and Essences}, 2013).
  Their mapping onto the admissibility schema is the authors'
  interpretive proposal.
\end{established}

\subsection{Cognition as analogy-driven reduction}

Hofstadter's central claim is that cognition is not primarily
rule-following or symbolic deduction but \emph{analogy-making}:
concepts are ``bundles of analogies'' and thought proceeds through
fluid transitions between them.

The process Hofstadter calls the \emph{central cognitive loop}
is structurally isomorphic to explicit substitution:
\begin{enumerate}
  \item a stored conceptual structure (chunk) is activated;
  \item partial unpacking into working memory occurs;
  \item analogical perception propagates through the unpacked structure;
  \item new conceptual activations emerge and are re-chunked.
\end{enumerate}

Compare this to $\beta$-reduction:
\begin{enumerate}
  \item a $\lambda$-abstraction is activated by application;
  \item the argument is substituted through the scope;
  \item local redexes are exposed and further reduced;
  \item a new (possibly partially reduced) term emerges.
\end{enumerate}

\subsection{Chunking as nested bubble formation}

Hofstadter's \emph{chunking} operation recursively compresses conceptual
structures into larger units.
Small perceptual features become concepts; concepts become
higher-order schema; schema become background assumptions that operate
below the threshold of attention.

In Spherepop terms: each chunk is a bubble whose interior has been
stably collapsed and whose membrane presents only a high-level
interface to further interactions.
Nested chunks are nested bubbles.
Recursive chunking is the construction of deeply nested bubble hierarchies
in which interior structure has been successively collapsed into stable
sub-topologies.

\begin{definition}[Cognitive admissibility]
  A perceptual input $\Delta$ is \emph{cognitively admissible}
  relative to a conceptual chunk $\mathbf{b}$ if $\Delta$ activates
  one or more attachment sites in $\partial(\mathbf{b})$ through
  analogical resonance, i.e., if a structural similarity mapping
  exists between $\Delta$ and components of $I(\mathbf{b})$.
\end{definition}

\subsection{Analogy, slippage, and fuzzy collapse}

Hofstadter emphasizes that concepts are \emph{blurry attractors},
not rigid symbolic containers.
Analogical mapping is never exact: ``slippage'' occurs when a mapping
activates a related but non-identical concept.

Standard $\lambda$-calculus is discrete and exact: substitution
either succeeds or fails.
Spherepop, by contrast, can in principle accommodate \emph{partial}
and \emph{graded} admissibility: a bubble whose interior is only
partially compatible with an incoming perturbation can undergo
a \emph{partial collapse} that leaves residual open sites.
This is precisely the formal structure required for
analogy-making, semantic slippage, contextual reinterpretation,
and metaphorical extension.

\begin{table}[h!]
\centering
\caption{Three-way correspondence: $\lambda$-calculus, Hofstadter cognition, Spherepop.}
\label{tab:threeway}
\begin{tabular}{lll}
\toprule
\textbf{$\lambda$-calculus} & \textbf{Hofstadter cognition} & \textbf{Spherepop} \\
\midrule
Variable              & Open conceptual slot        & Attachment site \\
Abstraction           & Encapsulated procedure      & Scoped bubble \\
Application           & Conceptual activation       & Bubble injection \\
$\beta$-reduction     & Analogical rebinding        & Membrane collapse \\
$\alpha$-conversion   & Context-sensitive reinterpretation & Historical relabeling \\
Normal form           & Stable conceptual equilibrium & Stable bubble topology \\
Divergence            & Recursive reminding loops   & Infinite bubbling \\
Confluence            & Semantic coherence          & Topological reconciliation \\
\bottomrule
\end{tabular}
\end{table}

\subsection{Communication and semantic transport}

Hofstadter describes language as \emph{compressed symbolic transport}:
a speaker encodes a complex meaning into a compact symbolic string;
the listener reconstructs meaning by ``adding water'' — activating
relevant conceptual structures and allowing them to partially re-expand.

In Spherepop: communication succeeds when the receiver's bubble
topology is sufficiently compatible with the sender's to admit a
successful membrane collapse.
Translation across languages or conceptual frameworks becomes
topology-preserving reduction across incompatible semantic manifolds,
and failure to communicate is a gluing failure in the sense of
Section~\ref{sec:unification}.

\subsection{Implication for Curry--Howard}

Under the Spherepop--Hofstadter synthesis, the Curry--Howard
correspondence acquires a cognitive reading:
\begin{itemize}
  \item \emph{types} are admissibility membranes (conceptual constraints);
  \item \emph{proofs} are stabilized analogy trajectories;
  \item \emph{programs} are constrained conceptual collapses;
  \item \emph{normalization} is cognitive coherence formation.
\end{itemize}
Cognition is a generalized proof-reduction process operating over
analogical topologies rather than purely syntactic strings.

% ================================================================
\section{Analogy as Reduction}
\label{sec:analogy-reduction}
% ================================================================

\begin{speculative}
  The argument in this section is a theoretical proposal synthesizing
  Johansson's linguistic framework with the Spherepop--Hofstadter
  synthesis developed in the preceding sections.
  Johansson's specific claim that ``analogy formations could be viewed
  as a process of reduction'' is from his published work on
  construction and reduction in language acquisition.
  The generalization to a cross-domain admissibility-reduction schema
  is the authors' proposal.
\end{speculative}

\subsection{Johansson's precedent}

The present framework extends earlier reduction-based accounts of
cognition and language by proposing that analogy itself is a reduction
operation.
This claim is not intended merely as metaphor.

Johansson's \emph{Construction as Reduction} argues that language
acquisition is not fundamentally constructive in the traditional sense
but rather a cumulative process of selective reduction from a
surrounding linguistic population.
Crucially, Johansson explicitly states that ``analogy formations could
be viewed as a process of reduction,'' since analogy progressively
eliminates irregularity and reduces variation across linguistic forms.

The present work generalizes this insight beyond linguistic evolution
into a broader operational framework spanning computation, cognition,
semantics, and physical systems.

\subsection{Analogy as primary reduction mechanism}

In ordinary symbolic models, analogy is treated as a secondary or
heuristic process operating over already-defined structures.
Under the present framework, analogy becomes a \emph{primary} reduction
mechanism through which structures themselves stabilize.

The essential claim:
cognition survives precisely because it continuously reduces the number
of admissible trajectories available to interpretation and action.
Analogical perception performs this reduction by collapsing distinct
configurations into families of structurally compatible transformations.

The analogical operator may be written:
\[
  \mathcal{A} : \Bcal_1 \rightsquigarrow \Bcal_2
\]
where $\mathcal{A}$ reduces the distance between admissible structural
configurations $\Bcal_1$ and $\Bcal_2$ by identifying a
constraint-preserving mapping between their interiors.

This reformulates analogy with formal precision:
\emph{two structures are analogically related if a constrained
transformation preserves enough invariant structure for coherent
rebinding}.
This is the admissibility-transport reading of analogy:
vastly more rigorous than simply ``finding similarities.''

\subsection{Cross-domain instances}

The same reduction grammar appears across all domains treated in this
monograph:
\begin{itemize}
  \item In \textbf{linguistic evolution}: analogy reduces morphological
        irregularity across a linguistic population;
  \item In \textbf{semantic cognition}: analogy reduces interpretive
        uncertainty by collapsing competing conceptual trajectories;
  \item In \textbf{RSVP field dynamics}: entropy gradients reduce
        admissible future states toward locally stable attractors;
  \item In \textbf{sheaf semantics}: successful gluing reduces local
        semantic incompatibilities into coherent global sections;
  \item In \textbf{Marine processing}: salience extraction reduces
        high-dimensional waveform variability into stable jitter
        structures;
  \item In \textbf{quantum interference}: coherent phase interaction
        reduces propagation histories into measurable amplitudes.
\end{itemize}

What recurs is not literal ontological identity but a shared
operational grammar: \emph{stable systems preserve admissibility
through local reduction}.

\subsection{Construction through recursive reduction}

This yields a reversal of the traditional opposition between creativity
and reduction.
Construction may emerge through recursive reduction processes operating
across populations, histories, and semantic topologies.
Johansson's formulation already hints at this inversion when language
formation is described as a ``stasis between constructive and
eliminative reduction.''

The present framework extends that principle: reduction is not merely
elimination.
Properly understood, reduction is the mechanism through which coherence
emerges from excess possibility.

This directly extends the monograph's central thesis:
\begin{quote}
  \textbf{Analogy is admissibility-preserving transport.
  $\beta$-reduction is its symbolic special case.
  Membrane collapse is its geometric realization.
  Thermodynamic relaxation is its physical realization.}
\end{quote}

% ================================================================
\part{Biological Instantiation: Multimodal Lineage Tracing}
% ================================================================

% ================================================================
\section{Cellular Differentiation as Admissibility-Constrained Collapse}
\label{sec:lineage}
% ================================================================

\begin{established}
  The biological material in this section is drawn from Wang, He, and
  Hu (2026), ``Computational approaches for multimodal lineage tracing,''
  \textit{Nature Reviews Genetics}, \doi{10.1038/s41576-026-00969-9}.
  Specific characterizations of lineage-tracing methodology, ancestral
  state reconstruction, and computational frameworks follow that review.
  The interpretation of this material through the admissibility framework
  is the authors' proposal and does not represent the claims of the
  original authors.
\end{established}

\subsection{The underdetermination of cellular identity by present state}

In classical single-cell biology, cellular identity is treated as a
point in a high-dimensional molecular state space, typically defined by a
transcriptomic profile specifying the expression level of each gene in the
cell's genome.
Two cells with similar transcriptomic profiles are treated as similar
cells; two cells with divergent profiles are treated as distinct.
This framework is ontologically static: it describes what a cell is at
the moment of measurement, not how it came to be what it is or what it
can become.

The review by Wang, He, and Hu identifies a fundamental limitation of this
static framework.
Transcriptomic similarity does not reliably imply common developmental
ancestry, because phenotypically similar cells may emerge through entirely
different developmental histories.
Conversely, cells that appear transcriptomically distinct may share recent
common progenitors and therefore occupy the same admissibility basin with
respect to future developmental potential.
The present configuration of a cell underdetermines its generative history,
its inherited constraint structure, and the space of its admissible futures.

This is a precise biological instance of the underconstraint problem
developed in Section~\ref{sec:underconstraint}.
A transcriptomic measurement generates a local section of the developmental
manifold at a single moment, but the space of developmental trajectories
compatible with that measurement is large.
Multiple incompatible developmental histories may produce indistinguishable
terminal states.
Conversely, the same ancestral state may produce divergent terminal states
depending on the history of environmental perturbations, signaling inputs,
and stochastic gene-expression fluctuations encountered along the way.
In the language of the admissibility framework, present configuration alone
does not determine admissibility history, and admissibility history is what
constrains admissible futures.

\subsection{Lineage tracing as historical admissibility reconstruction}

Multimodal lineage tracing addresses this limitation by coupling
transcriptomic measurements with orthogonal records of developmental
ancestry.
These lineage records take several forms, including somatic mutations
accumulated during cell division, CRISPR-induced barcodes introduced
at defined developmental stages, DNA methylation patterns inherited across
cell generations, and spatial position within tissue architecture.
Each modality encodes a different aspect of the cell's developmental
history, and their integration progressively constrains the space of
developmental trajectories compatible with the observed present state.

From the perspective of the admissibility framework, lineage tracing is
an attempt to reconstruct the sequence of admissibility regions
$\Bcal_0 \to \Bcal_1 \to \cdots \to \Bcal_n$ that generated the observed
terminal configuration $\Bcal_n$.
Each transition in this sequence is a membrane collapse in the sense of
Definition~\ref{def:R}: a bounded developmental state accepts a compatible
perturbation (a signaling input, a morphogen gradient, a mechanical force,
a stochastic transcriptional fluctuation), internally reorganizes its
constraint structure (chromatin remodeling, transcription factor
redistribution, metabolic reprogramming), and collapses into a new
admissible successor state while preserving certain invariants (cell
viability, genomic integrity, tissue compatibility).
The invariant that lineage tracing tracks is ancestral identity: the
inherited sequence of developmental collapses that produced the current
state.

The review characterizes this problem precisely when it describes
multimodal lineage tracing as attempting to reconstruct ``trajectory
inference'' and ``ancestral state estimation'' across the developmental
manifold.
The key challenge is that single-cell sequencing is destructive: the cell
is lysed to extract its molecular contents, so the measured cell no longer
exists after measurement.
Ancestral states are therefore fundamentally unobserved, and computational
systems must reconstruct them from the terminal-leaf observations alone,
inferring which historical sequences of admissible collapses could have
generated the observed distribution of present states.

\subsection{Disparate multimodality and constraint reconciliation}

A central technical challenge identified in the review is what the authors
call ``disparate multimodality'': lineage relationships are discrete and
tree-structured, while molecular profiles are continuous and
high-dimensional.
The lineage tree assigns each cell a discrete ancestry path through a
bifurcating graph, while the transcriptomic state space assigns each cell
a position in a continuous high-dimensional manifold.
These two representations encode fundamentally different aspects of the
cell's identity and are not immediately reconcilable by standard
mathematical tools.

This disparate multimodality is, in the language of the admissibility
framework, the problem of reconciling historical constraint topology with
instantaneous manifold geometry.
The lineage tree functions as a constrained admissibility skeleton: it
records which developmental collapses have occurred in what order, and
each branch point represents an irreversible commitment to a developmental
trajectory that closes off other potential futures.
The transcriptomic manifold is the continuous semantic field in which the
consequences of those discrete commitments are expressed.

Formally, let $T$ denote the lineage tree, with nodes representing cell
generations and edges representing division events, and let
$\mathcal{M}$ denote the transcriptomic manifold with points representing
cell states.
The reconstruction problem requires a map $\phi : T \to \mathcal{M}$ that
is consistent with both the discrete ancestral relationships encoded in $T$
and the continuous trajectory structure observed in $\mathcal{M}$.
Such a map exists globally if and only if the local section assignments
are pairwise compatible and satisfy the cocycle condition: precisely the
sheaf-gluing condition of Appendix~\ref{app:sheaf}.
When the gluing condition fails, the lineage reconstruction is obstructed,
and the system cannot produce a globally coherent developmental history
consistent with all measurements simultaneously.

\subsection{Developmental hallucination as gluing obstruction}

The parallel with semantic hallucination identified in Section~\ref{sec:unification}
and Appendix~\ref{app:sheaf} is not merely formal.
In both cases, local compatibility does not guarantee global coherence.

In the language model case, a system may produce locally fluent and
pairwise-consistent text while generating a globally incoherent factual
structure.
The failure is detected only by checking triple overlaps rather than
pairwise ones: the first \v{C}ech cohomology class $[\omega] \in
\check{H}^1(X, \mathcal{F})$ is non-zero.

In the developmental biology case, the analogous failure occurs when
locally plausible state assignments to individual cells cannot be assembled
into a globally admissible lineage reconstruction.
Two cells may each be individually consistent with their assigned
transcriptomic profiles, and they may even be pairwise-consistent in the
sense that their profiles are compatible with a common progenitor under a
simple differentiation model, but the full reconstruction across all
cells in the dataset fails because the implied ancestral trajectories
violate the constraints imposed by the lineage tree structure.

This is what may be called \emph{developmental hallucination}: a
computational system assigns locally plausible cell states and cell
transitions, but the resulting developmental history is globally
inadmissible.
It implies progenitor states that are biologically implausible, transition
sequences that violate chromatin accessibility constraints, or ancestry
assignments that contradict the discrete lineage topology.
The hallucination is invisible at the level of individual cell-type
assignments and becomes apparent only in the global reconstruction.

The review identifies optimal transport methods such as LineageOT and
moslin as computational strategies specifically designed to address this
problem.
These methods seek the most economical mapping between developmental states
under lineage constraints, minimizing admissibility violation while
preserving inherited structural relationships.
In the admissibility framework, these methods are approximating the
globally consistent section of the developmental sheaf: the unique
assignment of cells to developmental trajectories that satisfies both the
local transcriptomic compatibility conditions and the global ancestral
coherence requirements.

\subsection{Progressive potency restriction as admissibility contraction}

The review's concept of ``progressive potency restriction'' is one of
the most structurally informative formulations in the paper from the
perspective of the admissibility framework.
As cells differentiate, they progressively lose the capacity to adopt
alternative cell fates.
A totipotent zygote can in principle give rise to any cell type.
A pluripotent stem cell can give rise to most cell types.
A committed progenitor can give rise to a restricted set of related cell
types.
A terminally differentiated cell typically cannot give rise to any other
cell type at all.

This progression is not merely a statistical observation about which cell
types tend to appear at which developmental stages.
It is a constraint on the admissible future manifold of each cell.
As differentiation proceeds, the set of admissible successor states
contracts.
Each developmental collapse irreversibly eliminates trajectories from the
cell's future possibility space.
The contraction is not merely a reduction in the number of available gene
expression states but a reduction in the dimensionality of the admissibility
region itself: the constraints that govern which transitions are biologically
realizable become progressively tighter as inherited chromatin structure,
DNA methylation patterns, and committed transcription factor networks
accumulate.

In the formal language of Appendix~\ref{app:rsvp}, this is entropy reduction
in the specific sense of the RSVP framework: not thermodynamic entropy alone,
but the logarithmic volume of admissible future trajectories.
The entropy of an undifferentiated cell with respect to its developmental
future is high; the entropy of a terminally differentiated cell is zero or
near zero.
Differentiation is the progressive reduction of this developmental entropy
through irreversible admissibility collapse.

This interpretation is consistent with the review's observation that
computational approaches to developmental dynamics often model cell
differentiation through Waddington-landscape-style potential energy
surfaces, in which cells occupy attractors separated by potential barriers.
In the admissibility framework, these attractors are admissible equilibria
of the developmental admissibility functional, the barriers are
inadmissibility regions that require non-admissible perturbations to
traverse, and the landscape itself changes over developmental time as
the constraint structure inherited from prior collapses modifies the
admissibility topology.

\subsection{Deep learning as trajectory compression}

The review's treatment of deep-learning frameworks for multimodal lineage
analysis is philosophically significant because it illustrates how modern
computational systems operationalize the admissibility framework
implicitly.
Methods such as LineageVAE, PORCELAN, and the deep lineage frameworks
reviewed in the paper construct latent embeddings of cell state that are
simultaneously constrained by transcriptomic similarity and lineage
topology.

In the admissibility framework, these methods are learning compressed
representations of the developmental admissibility manifold: lower-dimensional
summaries that preserve the historical constraint structure required for
coherent trajectory inference.
The variational autoencoder objective in LineageVAE, for instance, balances
a reconstruction term (fidelity to observed transcriptomic state) against
a regularity term (prior structure on the latent space), while the lineage
coupling introduces an additional constraint that latent representations
of ancestrally related cells must lie in geometrically consistent regions
of the latent manifold.

This is formally analogous to the CLIO projection framework introduced
in Section~\ref{sec:rsvp-synthesis}.
In CLIO, the observer navigates a compressed projection $\pi : X \to M$
from a high-dimensional ontic trajectory space $X$ to an operational
measurement manifold $M$, preserving only the dynamically relevant
constraint relations.
The deep lineage methods are learning precisely this kind of projection:
a compression that preserves the admissibility-bearing constraint
structure of the developmental history while discarding the high-dimensional
details that are not required for trajectory inference.

The review notes that these methods remain limited by the heterogeneity
and incompleteness of multimodal lineage data, and by the difficulty of
learning latent representations that correctly balance the disparate
topological structures of lineage trees and continuous transcriptomic
manifolds.
In the admissibility framework, this limitation corresponds to the
difficulty of learning an admissible projection when the admissibility
structure of the full developmental manifold is accessible only through
sparse, destructive, and temporally incomplete measurements.
The deep-learning system is attempting to infer $\Rcal$ from observations
of $\Bcal_{i+1}$ alone, without access to the intermediate collapse
operators or the full history $H(\Bcal_{i+1})$.

\subsection{Admissibility as biological necessity, not metaphorical decoration}

The most important philosophical consequence of this biological grounding
is that admissibility is here not a speculative theoretical overlay but a
structural necessity of the system.
Cellular differentiation cannot explore arbitrary molecular configurations.
The space of biologically realizable developmental trajectories is a small
fraction of the space of all conceivable transitions between molecular
states, constrained by heredity, chromatin architecture, spatial geometry,
signaling network topology, energy budget, and the irreversibility of
committed developmental collapses.

Without admissibility constraints, multicellular coherence would dissolve.
If cells could transition arbitrarily between molecular states, the
coordination between cells required for tissue organization, organ
formation, and systemic physiology would become combinatorially impossible.
The organism survives precisely because developmental collapse is
constrained: each cell's trajectory is bounded by inherited admissibility
structure, and that structure propagates across cell generations through
mechanisms of epigenetic inheritance, gap junction communication,
morphogen gradients, and mechanical coupling.

The admissibility framework therefore finds in multimodal lineage tracing
not an analogy but an instantiation.
The admissibility regions are the developmentally accessible chromatin
configurations.
The perturbations are the signaling inputs and environmental fluctuations
to which the cell is exposed.
The compatibility conditions are the biochemical requirements for signal
transduction.
The collapse operators are the gene regulatory network dynamics that
integrate the perturbation and reorganize the cell's transcriptional state.
The invariant is cellular viability and genomic integrity.
The normal form is the terminally differentiated cell type.
The failure mode is developmental mis-specification, in which a cell
collapses into an inadmissible state, typically resulting in cell death,
oncogenic transformation, or developmental arrest.

The review by Wang, He, and Hu (2026) thus provides an empirically grounded,
computationally active, and formally tractable domain in which the
admissibility principle is not a philosophical hypothesis but a
computational necessity: the reconstruction of developmental history
requires admissibility constraints, and those constraints are
mathematically specifiable, biologically motivated, and progressively
recoverable through the integration of multiple heterogeneous data
modalities.

% ================================================================
\part{Established Physics}
% ================================================================

% ================================================================
\section{Quantum Weak-Value Dwell Times}
\label{sec:quantum}
% ================================================================

\begin{established}
  Results in this section are drawn from:
  APL Quantum \textbf{2}, 036108 (2025) on atomic excitation times
  for transmitted and scattered photons.
  Statements attributed to that paper reflect its conclusions;
  the admissibility mapping is the authors' interpretation.
\end{established}

\subsection{The photon dwell-time problem}

When a single photon propagates through a cloud of two-level atoms,
it may briefly exist as a collective atomic excitation
(a Dicke-like polariton state) before being re-emitted.
The question of how long the photon ``spends'' as an atomic excitation
has a physically non-trivial answer.

For \emph{transmitted} photons, the atomic excitation time
$\tau_T$ equals the net group delay:
\[
  \tau_T = \tau_{\mathrm{group}}.
\]
In resonant regimes where dispersion produces negative group delay
(the ``superluminal'' regime), $\tau_T$ is negative.
This is not a violation of causality but an instance of an
\emph{anomalous weak value}: a postselection effect arising from
quantum interference between two pathways.

For \emph{scattered} photons, the excitation time involves an
additional contribution:
\[
  \tau_S = \langle \tau_{\mathrm{group}} \rangle_{\mathrm{ens}}
           + \tau_{\mathrm{Wigner}},
\]
where $\tau_{\mathrm{Wigner}}$ is the Wigner time delay associated
with elastic scattering.

\subsection{Anomalous weak values as interference-mediated admissibility}

The negative-time result arises from interference between:
\begin{enumerate}
  \item a pathway in which the photon passes without excitation, and
  \item a pathway in which the photon converts to an atomic excitation
        and coherently transfers back into the photonic mode.
\end{enumerate}
Destructive interference between these pathways produces a net
negative pointer shift — an anomalous weak value.

\begin{note}[Caveat on ``negative time'']
  The phrase ``the photon spends a negative amount of time as an
  excitation'' is shorthand for a postselected weak-value measurement.
  It does not imply backward causation.
  The weak value is a property of the measurement protocol (pre- and
  post-selection), not a directly observable duration in the
  classical sense.
\end{note}

\subsection{Admissibility mapping}

The quantum dwell-time experiment instantiates the admissibility schema
as follows:
\begin{itemize}
  \item \textbf{Admissibility region} $\Bcal_i$: the atomic medium
        with its resonance structure and spectral response function;
  \item \textbf{Perturbation} $\Delta_i$: the incoming photon wavepacket;
  \item \textbf{Constraint / compatibility test}: the photon's spectral
        components must overlap the atomic resonance for excitation to
        be possible; the interference condition determines
        admissibility of the excitation pathway;
  \item \textbf{Collapse operator} $\Rcal$: destructive/constructive
        interference between pathways;
  \item \textbf{Stable output} $\Bcal_{i+1}$: transmitted photon with
        modified group delay, or scattered photon with Wigner delay.
  \item \textbf{Invariant preserved}: unitarity / probability conservation.
\end{itemize}

The confluence principle here is the Born rule: different measurement
protocols on the same quantum state must produce statistically consistent
results.

% ================================================================
\section{DC-OPA and Ultrafast Nonlinear Optics}
\label{sec:optics}
% ================================================================

\begin{established}
  Results in this section are drawn from the advanced DC-OPA
  experiment reported in \textit{Nature Photonics} (2024--25):
  a dual-chirped optical parametric amplification system using
  BiBO and MgO:LN crystals.
  Specific figures (pulse duration, peak power, CEP stability)
  are taken from that source.
  Claims about water-window harmonic generation efficiency and
  attosecond pulse duration are the authors' reading of the
  literature and should be independently verified.
\end{established}

\subsection{Dual-chirped optical parametric amplification}

High-order harmonic generation (HHG) for attosecond science requires
few-cycle mid-infrared (MIR) drivers.
Standard OPCPA (optical parametric chirped-pulse amplification)
stretches a seed pulse, amplifies it parametrically, and compresses
it back to short duration.

\emph{Dual-chirped} OPA (DC-OPA) goes further: both the pump
\emph{and} the seed are chirped.
This relaxes phase-matching bandwidth constraints and allows amplification
over more than one octave using heterogeneous nonlinear crystals
(BiBO for the first stage; MgO-doped lithium niobate for the second).

The reported performance of the advanced DC-OPA system:
\begin{itemize}
  \item Output energy: 53 mJ;
  \item Pulse duration: 8.58 fs at 2.44~$\mu$m central wavelength
        ($\approx 1.05$ optical cycles);
  \item Peak power: $\sim$6 TW;
  \item CEP stability: 228 mrad.
\end{itemize}

\subsection{Phase-matched energy redistribution as constrained collapse}

The physics of parametric amplification is governed by phase-matching:
energy transfer from pump to signal and idler is efficient only when
the wavevectors satisfy
\[
  \Delta k = k_p - k_s - k_i = 0.
\]
The nonlinear crystal is the constraint membrane;
the incoming pump and seed are the perturbation;
phase-matched energy redistribution is the collapse operation;
the output compressed pulse is the stable attractor.

\begin{itemize}
  \item \textbf{Admissibility region} $\Bcal_i$: the nonlinear crystal
        with its phase-matching bandwidth and group-velocity dispersion profile;
  \item \textbf{Perturbation} $\Delta_i$: the chirped pump and seed pulses;
  \item \textbf{Compatibility condition}: $\Delta k \approx 0$ over
        the amplification bandwidth;
  \item \textbf{Collapse operator} $\Rcal$: nonlinear three-wave mixing
        under phase-matching constraint;
  \item \textbf{Stable output} $\Bcal_{i+1}$: compressed MIR pulse
        with coherent phase structure;
  \item \textbf{Invariant}: photon number conservation (Manley--Rowe
        relations); CEP coherence.
\end{itemize}

\subsection{HHG and attosecond science}

The compressed MIR pulses drive HHG in gas targets.
An electron is ionized by the intense field, accelerated, and
recombines with the ion, emitting a burst of XUV or soft X-ray radiation.
The three-step model (ionization, propagation, recollision) is itself
an admissibility sequence: recollision is gated by the phase structure
of the driving field, and only electron trajectories that satisfy the
recollision condition contribute to harmonic emission.

The extension to isolated attosecond pulses (IAPs) requires spectral
selection from a supercontinuum region where contributions from multiple
half-cycles do not interfere destructively.
This is a second-level admissibility gate on top of the first.

% ================================================================
\part{Speculative and Prototypical Layers}
% ================================================================

% ================================================================
\section{RSVP: Thermodynamic Field Semantics}
\label{sec:rsvp}
% ================================================================

\begin{speculative}
  The Relativistic Scalar-Vector Plenum (RSVP) framework is a
  theoretical proposal under development by the author.
  Parameters such as the critical coupling $\lambda_c \approx 0.42$
  and the entropy production threshold $\dot{\Sigma}_{\mathrm{crit}}
  = 2.1\,\nats/\text{generation}$ are \emph{model-derived conjectural
  thresholds}, not empirically measured physical constants.
  They should not be read as established results of the same
  evidential status as the quantum-optics or ultrafast-laser material
  in Part~II.
\end{speculative}

\subsection{Framework overview}

RSVP models the evolution of physical and cognitive systems through
a coupled field theory over three interacting fields:
\begin{itemize}
  \item $\Phi$: a scalar entropy-potential field;
  \item $\mathbf{v}$: a vector flow field encoding directed processes;
  \item $S$: an entropy density field.
\end{itemize}
The central postulate is that admissible system trajectories are
those along which entropy production remains bounded:
$\dot{S} \geq 0$ locally (second law), but
$\dot{S} \leq \dot{\Sigma}_{\mathrm{crit}}$ globally for stable
operation (stability condition).

\subsection{RSVP as thermodynamic admissibility collapse}

In the RSVP framework, a local region of the field undergoes
admissibility-constrained collapse as follows:
\begin{itemize}
  \item \textbf{Admissibility region} $\Bcal_i$: a connected region
        of $(\Phi, \mathbf{v}, S)$ field space with locally defined
        entropy gradient $\nabla S$;
  \item \textbf{Perturbation} $\Delta_i$: an incoming entropy flux
        or directed flow perturbation;
  \item \textbf{Compatibility condition}: the perturbation must be
        absorbable without driving $\dot{S}$ beyond
        $\dot{\Sigma}_{\mathrm{crit}}$;
  \item \textbf{Collapse operator} $\Rcal$: entropy-gradient-constrained
        field relaxation;
  \item \textbf{Stable output} $\Bcal_{i+1}$: a new field configuration
        with lower free energy and admissible entropy production;
  \item \textbf{Invariant}: total entropy non-decrease; field continuity.
\end{itemize}

\subsection{The thermodynamic reading of $\beta$-normal form}

In $\lambda$-calculus, a $\beta$-normal form is a term with no
remaining redexes: computation has reached a stable configuration.
In RSVP, the corresponding notion is a field state in which no
admissible entropy-decreasing trajectory remains locally available:
a thermodynamic equilibrium or metastable attractor.

The analogy is not that computation and thermodynamics are the same
physics, but that both are instances of the generalized reduction
schema in which ``no further collapse is possible'' constitutes the
normal form.

\subsection{RSVP parameters as scaling hypotheses}

The numerical thresholds $\lambda_c \approx 0.42$ and
$\dot{\Sigma}_{\mathrm{crit}} = 2.1\,\nats/\mathrm{generation}$
are \emph{provisional scaling hypotheses} defining admissible behavioral
regimes within RSVP simulation work.
They are not empirical constants, not derived from first principles
within an established physical theory, and not published in
peer-reviewed venues.

The appropriate analogy is Reynolds-number thresholds in early
turbulence research: numbers that organized experimental intuition
and defined qualitative regime boundaries before quantitative derivation
was available.
The RSVP constants should be treated as heuristic bifurcation indicators
subject to revision as the framework is developed.

\subsection{Thermodynamic constraints versus impossibility proofs}
\label{sec:foom-caveat}

The RSVP framework motivates a thermodynamic argument against
unbounded recursive self-improvement (``FOOM''):
\begin{enumerate}
  \item Any cognitive subsystem increasing organizational complexity
        must export entropy to its environment
        (second law applied to open systems).
  \item The entropy export rate is bounded by thermal conductivity
        and effective surface area.
  \item Therefore, the rate of self-improvement is bounded by a
        finite thermodynamic flux.
\end{enumerate}
This argument is a \emph{scaling constraint}, not a hard impossibility
theorem.
Several mechanisms could alter the bounds:
\begin{itemize}
  \item \emph{Reversible computing}: Landauer's principle permits
        logically reversible operations to approach zero heat
        dissipation per step; a sufficiently reversible architecture
        could reduce entropy export per operation arbitrarily.
  \item \emph{Quantum adiabatic evolution}: adiabatic processes
        produce no entropy in the ideal limit.
  \item \emph{Distributed extension}: a system growing spatially
        can increase effective surface area, relaxing the area bound.
\end{itemize}
The correct claim is: \emph{classical, locally bounded self-improvement
faces finite thermodynamic constraints under the RSVP parameters as
currently defined}.
Whether those constraints bind depends on architecture and regime.
This is a philosophical motivation and a modeling choice, not a theorem
of the status of the second law itself.

% ================================================================
\section{Marine, Phoenix 8a, and MEM$|$8}
\label{sec:phoenix}
% ================================================================

\begin{prototype}
  Phoenix 8a, Marine pattern memory, and MEM$|$8 are engineering
  architectures at the prototype or proposal stage.
  Benchmarks quoted (RMS match accuracy $> 99.5\%$, voice cloning
  from minimal data) are internal test results that have not been
  independently replicated or published in a peer-reviewed venue.
  Claims about ``infinite frequency resolution'' and ``zero latency''
  are significant overstatements of information-theoretic limits
  and must be qualified (see Section~\ref{sec:marine-caveat}).
  The material here is presented as motivated engineering hypothesis,
  not established science.
\end{prototype}

\subsection{The core design philosophy}

Traditional digital audio processing converts continuous waveforms
into sampled pulse-code modulation (PCM) and analyzes them via the
Short-Time Fourier Transform (STFT) or Mel spectrogram.
The fundamental trade-off is the time-frequency uncertainty relation:
\[
  \Delta t \cdot \Delta f \;\geq\; \frac{1}{4\pi}.
\]
Large analysis windows provide high frequency resolution but poor
time resolution, destroying attack transients and micro-temporal
features that carry perceptual information about articulation,
emotional valence, and speaker identity.

The Marine engine proposes an alternative: \emph{time-domain interval
counting}.
Rather than sampling waveform amplitude at a fixed rate and
Fourier-transforming, Marine tracks the intervals between
zero-crossings and local extrema, extracting pitch, ADSR envelopes,
vibrato trajectories, and jitter signatures directly from the
temporal structure of the waveform.

\subsection{Admissibility mapping}

\begin{itemize}
  \item \textbf{Admissibility region} $\Bcal_i$: the temporal-interval
        structure of a waveform segment, interpreted as a constraint
        membrane over the space of admissible harmonic families;
  \item \textbf{Perturbation} $\Delta_i$: an incoming waveform event
        (zero-crossing, peak, onset);
  \item \textbf{Compatibility condition}: the event's temporal position
        must be consistent with the current harmonic-family hypothesis;
  \item \textbf{Collapse operator} $\Rcal$: jitter/salience extraction,
        ADSR envelope estimation, harmonic-lock detection;
  \item \textbf{Stable output} $\Bcal_{i+1}$: a compressed representation
        as a sparse set of harmonic families with associated envelope
        and jitter parameters;
  \item \textbf{Invariant}: temporal ordering; pitch contour continuity.
\end{itemize}

\subsection{Caveat: time-frequency uncertainty}
\label{sec:marine-caveat}

The phrase ``infinite frequency resolution with zero latency'' cannot
be taken literally: it violates the time-frequency uncertainty
principle, which is not an artifact of FFT implementation but a
consequence of the Fourier duality between time and frequency
representations.
What the Marine approach \emph{can} legitimately claim is:
\begin{itemize}
  \item resolution of \emph{pitch} from temporal periodicity without
        a fixed window function;
  \item lower latency than STFT for onset detection;
  \item extraction of temporal features (jitter, vibrato rate) that
        STFT discards by design.
\end{itemize}
These are meaningful advantages that should be claimed on their own
terms rather than through a physically incoherent slogan.

\subsection{MEM$|$8 as persistent attractor topology}

MEM$|$8 proposes to store compressed audio ``atoms'' (harmonic families
and residual burst structures) as activation seeds for a semantic
memory system.
The design goal is that memory retrieval proceeds by \emph{analogical
resonance} rather than hash-lookup: a query activates the stored atoms
whose structure is most compatible with the query's temporal-harmonic profile.

In Hofstadter's terms, this is a formal implementation of
analogy-driven memory reconstruction.
In Spherepop terms, stored atoms are bubble-topology snapshots;
retrieval is admissibility-gated membrane activation.
Whether current implementations of MEM$|$8 realize this design goal
at the claimed fidelity levels is a question that requires
independent benchmark evaluation.

% ================================================================
\section{The Generalized Reduction Operator}
\label{sec:unification}
% ================================================================

\subsection{Formal statement}

We now bring together the preceding analyses into a single schema.

\begin{definition}[Generalized reduction operator]
\label{def:R}
  Let $\Bcal$ be a class of \emph{bounded admissibility regions} and
  $\Dcal$ a class of \emph{perturbations}.
  A \emph{generalized reduction operator} is a partial function
  \[
    \Rcal \;:\; \Bcal \times \Dcal \;\rightharpoonup\; \Bcal
  \]
  satisfying:
  \begin{enumerate}[label=(\roman*)]
    \item \textbf{Admissibility gate}: $\Rcal(\Bcal_i, \Delta_i)$ is
          defined only when $\Delta_i$ is compatible with the constraint
          structure $\partial(\Bcal_i)$;
    \item \textbf{Invariant preservation}: there exists an invariant
          function $\iota : \Bcal \to I$ such that
          $\iota(\Rcal(\Bcal_i, \Delta_i)) = \iota(\Bcal_i)$
          whenever $\Rcal$ is defined;
    \item \textbf{Irreversibility}: the history $H$ is strictly extended
          by each application of $\Rcal$;
    \item \textbf{Locality}: $\Rcal(\Bcal_i, \Delta_i)$ depends only
          on the local constraint structure $\partial(\Bcal_i)$ and the
          interior $I(\Bcal_i)$, not on the global state of the system.
  \end{enumerate}
\end{definition}

\subsection{Necessary and sufficient conditions for instantiating
  $\Rcal$: what does \emph{not} qualify}

The reviewer of any draft of this work will correctly ask:
\emph{what would not count as an instance of $\Rcal$}?
Without a nontrivial exclusion criterion, the schema risks vacuity —
any constrained process with inputs and outputs could be
force-fit into the table.
This subsection addresses that objection directly.

\paragraph{Three necessary conditions beyond Definition~\ref{def:R}.}

A process qualifies as an admissibility-collapse system only if
all four conditions of Definition~\ref{def:R} hold
\emph{and} three additional structural requirements are met:

\begin{enumerate}
  \item \textbf{Recursive state-space restructuring.}
        Admissibility must actively change the future trajectory space,
        not merely filter outputs after the fact.
        A thermostat does not qualify: it enforces a temperature band
        but does not restructure the space of admissible future states.
        A lookup table does not qualify: it maps inputs to outputs
        without recursively reorganizing the constraint structure.
        The test is whether removing the admissibility gate would
        allow the system to reach qualitatively different macrostates.

  \item \textbf{Transportable invariant.}
        Collapse must propagate a coherent invariant through the
        transformation, not merely destroy structure.
        A shattered glass is not performing admissibility-guided
        reduction: no coherent invariant ($\iota$) survives the
        transformation in a form that constrains future evolution.
        The invariant must be \emph{transportable} --- usable by
        downstream processes as a constraint.

  \item \textbf{Path-sensitive historical dependence.}
        The system's admissibility structure must depend on its
        reduction history, not merely its current state.
        This is condition~(iii) of Definition~\ref{def:R} made
        operationally precise.
        A Markov process with no memory of prior states does not
        qualify unless its transition kernel is itself modified by
        prior reductions.
\end{enumerate}

\paragraph{Canonical non-instances.}

The following processes explicitly do not instantiate $\Rcal$:
\begin{itemize}
  \item \emph{Thermostats and bang-bang controllers}:
        constraint without state-space restructuring.
  \item \emph{Hash lookups and lookup tables}:
        output filtering without invariant propagation.
  \item \emph{Markovian random walks}:
        locally defined transitions without historical dependence.
  \item \emph{Physical destruction} (shattering, combustion without
        product reuse): transformation without transportable invariant.
  \item \emph{Arbitrary compression} (e.g., lossy JPEG):
        projection without admissibility gate or invariant preservation.
\end{itemize}

\paragraph{Why biological development qualifies.}

A fertilized egg undergoing development is a canonical
positive instance.
Each admissible developmental collapse (cell fate determination,
tissue boundary formation, organogenesis) recursively restricts
future developmental possibilities while preserving organismic
continuity.
The space of possible futures narrows through a sequence of
history-dependent admissibility collapses, each propagating
structural invariants (positional information, morphogen gradients,
lineage identity) that constrain subsequent stages.
This is precisely the intended meaning of $\Rcal$.

\subsection{Domain specializations}

Table~\ref{tab:atlas} summarizes how each domain instantiates
Definition~\ref{def:R}.

\begin{longtable}{p{2.8cm}p{3cm}p{3cm}p{3cm}p{3cm}}
\caption{Concept atlas: domains as specializations of $\Rcal$.}
\label{tab:atlas}\\
\toprule
\textbf{Domain} & \textbf{Admissibility region $\Bcal_i$} &
\textbf{Perturbation $\Delta_i$} & \textbf{Collapse operator $\Rcal$} &
\textbf{Invariant / failure mode} \\
\midrule
\endfirsthead
\multicolumn{5}{l}{\small\itshape Table \ref{tab:atlas} continued.}\\
\toprule
\textbf{Domain} & \textbf{Admissibility region} &
\textbf{Perturbation} & \textbf{Collapse operator} &
\textbf{Invariant / failure mode} \\
\midrule
\endhead
\midrule
\multicolumn{5}{r}{\small\itshape continued on next page}\\
\endfoot
\bottomrule
\endlastfoot
$\lambda$-calculus &
  Scoped abstraction $\Lam{x}{M}$ &
  Argument $N$ &
  $\beta$-reduction: $\subst{M}{x}{N}$ &
  $\beta$-normal form / divergence ($\Omega$) \\[4pt]

Spherepop &
  Scoped bubble with membrane $\partial$ &
  Compatible bubble $\mathbf{b}'$ &
  Membrane collapse $\oplus$ &
  Stable topology / infinite bubbling \\[4pt]

Hofstadter cognition &
  Conceptual chunk (bundle of analogies) &
  Perceptual input or memory cue &
  Analogical rebinding (central cognitive loop) &
  Stable concept / recursive reminding loop \\[4pt]

Quantum dwell time &
  Atomic medium (resonance + dispersion) &
  Photon wavepacket &
  Interference-mediated pathway collapse &
  Group delay (possibly anomalous weak value) / decoherence \\[4pt]

DC-OPA / HHG &
  Nonlinear crystal (phase-matching envelope) &
  Chirped pump + seed pulses &
  Three-wave mixing under $\Delta k = 0$ &
  Compressed CEP-stable pulse / group-velocity mismatch \\[4pt]

RSVP &
  $(\Phi, \mathbf{v}, S)$ field region &
  Entropy flux $\delta S$ &
  Gradient-constrained field relaxation &
  Thermodynamic attractor / supercritical runaway \\[4pt]

Marine / Phoenix 8a &
  Temporal-interval structure of waveform &
  Onset event (zero-crossing, peak) &
  Jitter + ADSR extraction, harmonic-lock &
  Sparse harmonic family representation / lock failure \\[4pt]

Sheaf semantics &
  Open cover of semantic patches &
  Overlap data between patches &
  Gluing map &
  Global section (coherent meaning) / gluing obstruction \\
\end{longtable}

\subsection{Why confluence appears across domains: family resemblance,
  not a universal theorem}

\begin{note}[Terminological precision]
  The phrase ``master convergence principle'' used in earlier drafts
  of this work is \emph{heuristic terminology}, not a proved universal
  theorem.
  The Church--Rosser theorem is a precisely stated and proved result
  about $\lambda$-calculus.
  The Born-rule consistency condition is a physical postulate about
  quantum measurement.
  The existence of global sheaf sections is a categorical result in
  algebraic topology.
  These are mathematically distinct statements that cannot be collapsed
  into a single theorem without a category-theoretic framework
  (e.g., showing that all are instances of limits preserving coherence
  in a 2-category) that is not provided here.
\end{note}

What can be legitimately claimed is that these domains exhibit a
\emph{family resemblance} generated by a common structural motif:
\emph{local compatibility reconciliation producing globally coherent
output}.

Consider the analogical reading of each convergence result:
\begin{itemize}
  \item In $\lambda$-calculus: incompatible substitution orders are
        reconciled through normalization to a common normal form.
  \item In sheaf theory: overlapping local descriptions are reconciled
        through gluing maps into a global section.
  \item In turbulence: incompatible local energy-transfer rates are
        reconciled through admissible dissipation into a single
        physically correct trajectory.
  \item In game theory: incompatible strategy profiles are reconciled
        through iterated dominance elimination into a common solution
        concept.
  \item In cognition: competing analogical interpretations are
        reconciled through contextual rebinding into stable conceptual
        equilibria.
\end{itemize}

The common element is not a shared theorem but a shared \emph{question}:
when a system faces locally incompatible pressures, does it possess a
coherent collapse pathway?
For $\lambda$-calculus the answer is the Church--Rosser theorem.
For the other domains, confluence is a hypothesis, an engineering
desideratum, or a phenomenon to be explained.

The admissibility principle is the claim that this question has the
same \emph{form} across domains, not that it has the same
\emph{answer}.
Recognizing the shared form is the contribution; proving or disproving
confluence in each domain is the research program.

\subsection{Sheaf obstruction and hallucination}

An important failure mode of admissibility-constrained collapse is
\emph{gluing obstruction}: local patches are each internally consistent
but cannot be assembled into a globally coherent state.

In sheaf theory: a presheaf satisfies local compatibility but fails
to admit a global section; the obstruction is measured by the first
\v{C}ech cohomology group $\check{H}^1$.

In cognitive terms: individual analogical fragments activate without
contradiction locally, but their combination produces a globally
incoherent representation.
This is a formal model of semantic hallucination in language models:
the system generates locally plausible token sequences whose global
factual structure violates world-knowledge constraints.

\begin{proposition}[Hallucination as gluing failure]
  Let $\{U_i\}$ be an open cover of a semantic space, and let
  $s_i \in \mathcal{F}(U_i)$ be locally consistent semantic patches.
  A \emph{hallucination} occurs when the cocycle condition
  $s_i|_{U_i \cap U_j} = s_j|_{U_i \cap U_j}$ fails for some $i, j$,
  so that no global section $s \in \mathcal{F}(X)$ restricting to
  all $s_i$ exists.
\end{proposition}

% ================================================================
\part{Convergent Literature and the Five-Layer Architecture}
% ================================================================

% ================================================================
\section{Game-Theoretic Admissibility: Recursive Elimination Under
  Common Knowledge}
\label{sec:game}
% ================================================================

\begin{established}
  The material in this section follows Bicchieri and Schulte (1996),
  ``Common Reasoning About Admissibility,''
  \textit{Erkenntnis} (revised 1996).
  The mapping onto the Spherepop and sheaf frameworks is the
  authors' interpretive proposal.
\end{established}

\subsection{Admissibility as a dynamical operator on possibility spaces}

Bicchieri and Schulte introduce a dimension of admissibility absent from
purely structural accounts: admissibility is not a static property of
states but a \emph{recursive elimination} process operating under
common knowledge constraints.

Their central result concerns strategic games.
A strategy $s_i$ for player $i$ is \emph{(weakly) admissible} if it is
not weakly dominated: no alternative strategy yields a better outcome
against every possible opponent strategy.
\emph{Common reasoning about admissibility} is the iterated process
by which players eliminate inadmissible strategies under the assumption
that all players are doing likewise.

\begin{definition}[Sequential proper admissibility]
  A strategy $s_i$ is \emph{sequentially properly admissible} in a game
  tree $T$ if $s_i$ is admissible at every information set $I_i$
  consistent with $s_i$, i.e., at every node reachable under $s_i$.
\end{definition}

The main theorem of Bicchieri and Schulte establishes that
in finite extensive games with perfect recall,
the strategies consistent with common reasoning about sequential proper
admissibility in the extensive form coincide exactly with those
consistent with common reasoning about admissibility in the strategic
form.
In other words: the solution given by recursive admissibility elimination
is representation-independent.

\subsection{Admissibility as trajectory pruning}

The operational schema of Bicchieri--Schulte is:
\[
  \text{possible trajectories}
  \;\longrightarrow\;
  \text{eliminate inadmissible branches}
  \;\longrightarrow\;
  \text{recurse under common knowledge}
  \;\longrightarrow\;
  \text{stable surviving trajectories.}
\]
This maps directly onto our generalized reduction operator:

\begin{itemize}
  \item \textbf{Admissibility region} $\Bcal_i$: the information set
        of player $i$, together with the probability assessment over
        opponent strategies;
  \item \textbf{Perturbation} $\Delta_i$: a new strategy or
        deviation possibility;
  \item \textbf{Compatibility condition}: the strategy must survive
        iterated weak-dominance elimination under strict coherence
        (every opponent strategy receives positive probability);
  \item \textbf{Collapse operator} $\Rcal$: iterated elimination of
        weakly dominated strategies under common knowledge;
  \item \textbf{Stable output}: the set of rationalizable, admissibility-
        consistent strategy profiles;
  \item \textbf{Invariant}: representation-independence
        (strategic-form and extensive-form admissibility coincide).
\end{itemize}

The key philosophical point is that admissibility here is not static.
It propagates globally through nested epistemic embeddings.
This is structurally identical to:
\begin{itemize}
  \item recursive sheaf gluing across overlapping open sets;
  \item Hofstadter's strange loops (cognition reasoning about its own
        reasoning);
  \item RSVP/CLIO recursive coherence filtering.
\end{itemize}

\begin{table}[h!]
\centering
\caption{Correspondence: game-theoretic admissibility and Spherepop.}
\label{tab:game-spherepop}
\begin{tabular}{ll}
\toprule
\textbf{Game theory (Bicchieri--Schulte)} & \textbf{Spherepop} \\
\midrule
Strategy                      & Bubble trajectory \\
Information set               & Local membrane state \\
Weak dominance                & Inadmissible collapse path \\
Iterated admissibility (IWD)  & Recursive membrane pruning \\
Sequential admissibility      & Path-consistent topology \\
Common knowledge of rationality & Shared admissibility manifold \\
Backward induction            & Bottom-up bubble resolution \\
Forward induction             & Top-down admissibility propagation \\
\bottomrule
\end{tabular}
\end{table}

\subsection{Backward and forward induction as admissibility projections}

Bicchieri and Schulte show that backward and forward induction
principles, often treated as mutually inconsistent, both emerge from
one underlying principle: common reasoning about admissibility.
Their unification parallels the monograph's broader claim that
apparently distinct reasoning paradigms (symbolic, topological,
thermodynamic, cognitive) are specializations of a common schema.

This also supplies the epistemic layer that earlier sections lacked.
RSVP provides thermodynamic admissibility.
Spherepop provides topological admissibility.
$\lambda$-calculus provides syntactic admissibility.
Bicchieri--Schulte provides \emph{epistemic admissibility under
recursive mutual constraint propagation}.
The four layers form a complete stack.

% ================================================================
\section{Turbulence Admissibility: Physical Realization Requires
  Selection Beyond Dynamics}
\label{sec:turbulence}
% ================================================================

\begin{established}
  Results in this section follow Glimm, Cheng, Sharp, and Kaman,
  ``A crisis for the verification and validation of turbulence
  simulations,'' published preprint (Los Alamos LA-UR-19-20285).
  The three-algorithm comparison and maximum-dissipation validation
  evidence are as reported there.
  The admissibility mapping is the authors' interpretation.
\end{established}

\subsection{The nonuniqueness problem in fluid dynamics}

The Euler equations for incompressible fluid flow,
\[
  \partial_t \mathbf{v} + (\mathbf{v}\cdot\nabla)\mathbf{v}
  = -\nabla p, \quad \nabla\cdot\mathbf{v} = 0,
\]
admit multiple physically distinct weak solutions even for identical
initial conditions.
This is not a numerical artifact: de Lellis and Székelyhidi have
constructed infinitely many energy-conserving wild solutions.
An additional physical principle is therefore required to select the
physically realized trajectory among mathematically admissible ones.
Glimm et al.\ call this the \emph{admissibility principle for
fluid dynamics}.

\subsection{Three admissibility criteria and their consequences}

Three physically motivated selection principles produce markedly
different simulations of Rayleigh--Taylor (RT) turbulent mixing:
\begin{enumerate}
  \item \textbf{Maximum dissipation rate} (FronTier algorithm):
        the local energy dissipation rate is maximized consistent with
        turbulent scaling laws.
  \item \textbf{Limited dissipation} (ILES / Miranda):
        the Reynolds stress is treated as a Gibbs-type numerical
        artifact to be minimized.
  \item \textbf{Near-zero dissipation} (MDNS):
        sub-grid-scale terms are suppressed based on a
        globally averaged Reynolds number.
\end{enumerate}
The three algorithms produce RT mixing rates $\alpha_b$ that differ
by factors of 2--3.
Validation against experimental data from Smeeton--Youngs supports the
maximum-dissipation principle.

The critical insight for the present framework:
\[
  \text{governing equations} + \text{admissibility principle}
  \;\longrightarrow\;
  \text{physically realized solution.}
\]
Dynamics alone are insufficient.

\subsection{Maximum entropy production as admissibility selector}

The paper identifies maximum entropy production as the admissibility
criterion consistent with experiment:
\[
  \dot{S} = \dot{S}_{\max}(x, t)
\]
locally, under the weak quasi-stationarity hypothesis.
This is a physical analogue of RSVP's thermodynamic admissibility:
in both cases, entropy production is not merely a consequence of
dynamics but a \emph{constraint selecting} which among many dynamically
consistent trajectories is realized.

\begin{note}[Reynolds stress as admissibility-bearing structure]
  The Reynolds stress $\overline{uv} - \bar{u}\bar{v}$ is not
  numerical noise.
  It encodes genuine sub-grid-scale admissibility information that
  cannot be projected away without altering the macroscopic trajectory.
  Discarding it (ILES/MDNS) produces physically incorrect mixing rates.
  This is a turbulence analogue of the semantic compression argument:
  projections that discard constraint-carrying structure alter
  admissibility and produce incorrect global behavior.
\end{note}

\subsection{Admissibility mapping for turbulence}

\begin{itemize}
  \item \textbf{Admissibility region} $\Bcal_i$: the flow field at
        resolved scales with its local turbulent statistics;
  \item \textbf{Perturbation} $\Delta_i$: sub-grid-scale Reynolds
        stress contributions;
  \item \textbf{Compatibility condition}: energy transfer must satisfy
        locally maximal dissipation under turbulent scaling laws;
  \item \textbf{Collapse operator} $\Rcal$: dynamic sub-grid-scale
        model (FronTier-style) realizing maximum local dissipation;
  \item \textbf{Stable output}: physically correct RT mixing rate
        $\alpha_b \approx 0.06$;
  \item \textbf{Invariant}: energy conservation (Manley--Rowe at the
        fluid level); Kolmogorov scaling.
  \item \textbf{Failure mode}: inadmissible suppression of Reynolds
        stress $\to$ $\alpha_b \approx 0.02$--$0.03$, contradicting
        experiment.
\end{itemize}

% ================================================================
\section{Observer Admissibility in Relativistic Physics}
\label{sec:relativity}
% ================================================================

\begin{speculative}
  The Modgil--Patil (2026) paper on observer admissibility and Planck
  invariance is not yet peer-reviewed in its current form and contains
  acknowledged mathematical gaps (particularly in the covariant
  formulation of $\gamma_g$ and the phenomenological global-bound
  formula).
  The Evangelista (preprint) and Zenteno (Zenodo, 2026) papers are
  similarly pre-publication.
  The material in this section is presented as a convergent theoretical
  program, not established physics.
  The mapping onto the admissibility schema is the authors' synthesis.
\end{speculative}

\subsection{The observerhood restriction strategy}

A cluster of recent proposals converges on a common move: rather than
modifying the dynamical laws of relativity to accommodate Planck-scale
physics, restrict the class of physically realizable observers.

The key thesis, stated in Modgil--Patil, is:
\begin{quote}
  Planck-scale limits can coexist with exact local Lorentz invariance
  when they are interpreted as operational bounds on admissible
  observers, rather than as indicators of symmetry breaking or
  fundamental spacetime discreteness.
\end{quote}

An observer is locally admissible when
\[
  g_{\mu\nu}V^\mu V^\nu > 0,
\]
i.e., the observer remains timelike relative to the local metric.
Global admissibility adds cosmological constraints:
\[
  \Delta t \leq T, \quad \Delta s \leq L, \quad m \leq M,
\]
where $T, L, M$ are finite characteristic scales of the universe.
The physically realizable domain is the intersection:
\[
  \mathcal{A} = \mathcal{A}_{\mathrm{local}} \cap \mathcal{A}_{\mathrm{global}}.
\]

\subsection{Admissibility mapping for relativistic observerhood}

\begin{itemize}
  \item \textbf{Admissibility region} $\Bcal_i$: the local spacetime
        neighborhood, together with global cosmological scale
        constraints;
  \item \textbf{Perturbation} $\Delta_i$: a candidate measurement
        protocol (localization attempt, boost, long-duration process);
  \item \textbf{Compatibility condition}:
        $g_{\mu\nu}V^\mu V^\nu > 0$ locally;
        $\Delta t \leq T$, $\Delta s \geq \ell_P$ globally;
  \item \textbf{Collapse operator} $\Rcal$: kinematic phase-boundary
        crossing (horizon encounter, Planck-scale localization
        attempt);
  \item \textbf{Stable output}: measurement result for an admissible
        observer; black-hole thermodynamics recovered without Lorentz
        violation;
  \item \textbf{Invariant}: exact local Lorentz symmetry; causal
        structure;
  \item \textbf{Failure mode}: inadmissible observer --- stationary
        observer at horizon, trans-Planckian localization attempt,
        infinite boost.
\end{itemize}

\subsection{The Kretschmann scalar as minimal admissibility detector}

Evangelista (preprint) provides a rigorous geometric instantiation
of the admissibility idea.
He imposes three requirements on an admissibility obstruction
functional $\mathcal{K}_\mathcal{A}[\Phi]$:
\begin{enumerate}
  \item local and covariant;
  \item scalar (observer-independent);
  \item non-degenerate in Ricci-flat (vacuum) spacetimes.
\end{enumerate}
The third condition is decisive: in Ricci-flat spacetimes,
$R = 0$ and $R_{\mu\nu}R^{\mu\nu} = 0$, so lower-order scalars vanish.
The \emph{Kretschmann scalar}
\[
  K = R_{\mu\nu\rho\sigma}R^{\mu\nu\rho\sigma}
\]
is the minimal quadratic invariant satisfying all three requirements:
it detects purely tidal (Weyl) curvature even in vacuum.

Admissibility may then be expressed as the curvature bound
\[
  \hat{K} = K\ell_o^4 \leq 1,
\]
which excludes curvature singularities from the physically realizable
domain.

\begin{remark}[Kretschmann as non-degenerate admissibility detector]
  The Kretschmann scalar is a physical instance of a general principle
  that recurs throughout this monograph:
  \emph{admissibility operators must remain sensitive to latent
  instability modes invisible to coarse observables}.
  In general relativity, lower-order Ricci scalars vanish in vacuum
  while dangerous Weyl curvature persists.
  Analogously, in language models, syntactic coherence fails to detect
  semantic hallucination; in ILES turbulence simulations, averaged
  Reynolds numbers fail to detect locally dangerous energy
  concentrations.
  The Kretschmann scalar formalizes what a non-degenerate admissibility
  detector must look like.
\end{remark}

\subsection{Zenteno: admissibility as dynamical stability criterion}

Zenteno (2026) formalizes admissibility through continuous penalty
functionals relative to constraint manifolds.
A structure is admissible if it exhibits \emph{bounded, stable, and
resolvable} behavior under soft constraint dynamics.
This generalizes classical existence-and-consistency criteria to
include dynamical persistence:

\begin{definition}[Zenteno admissibility]
  A mathematical structure $\mathcal{S}$ is \emph{admissible} relative
  to a constraint manifold $\mathcal{C}$ if the penalty functional
  \[
    \mathcal{P}[\mathcal{S}; \mathcal{C}] < \infty
  \]
  and the trajectory $\mathcal{S}(t)$ under constraint dynamics remains
  within a bounded neighborhood of $\mathcal{C}$ for all $t$.
\end{definition}

This framework naturally explains singularity exclusion:
singular geometries violate the boundedness condition rather than the
existence or consistency conditions.
The Paton (2026) pre-explanatory admissibility layer makes the same
move at the level of scientific modelling:
a system state must satisfy structural membership conditions before
any explanatory dynamics can meaningfully apply.

% ================================================================
\section{The Five-Layer Architecture}
\label{sec:fivelayers}
% ================================================================

The preceding sections have accumulated sufficient material to
articulate the monograph's full epistemic architecture explicitly.
We formalize five layers, each building on the previous:

\begin{description}
  \item[\textbf{Layer~I: Established empirical systems.}]
    Quantum optics (weak-value dwell times, APL Quantum 2025),
    ultrafast photonics (DC-OPA, Nature Photonics),
    gauge theory (standard),
    Rayleigh--Taylor turbulence validation (Glimm et al.).
    These are peer-reviewed results with experimental support.
    The admissibility schema maps onto them but is not required
    to explain them; they serve as \emph{test cases} for the
    schema's descriptive coherence.

  \item[\textbf{Layer~II: Engineering architectures and computational
    hypotheses.}]
    Marine temporal-interval analysis, Phoenix 8a, MEM$|$8,
    wave-memory semantic retrieval.
    These are prototype-stage proposals requiring independent
    benchmark validation.
    Their connection to the admissibility schema is motivational,
    not evidentiary.

  \item[\textbf{Layer~III: Formal abstraction layer.}]
    $\beta$-reduction, sheaf gluing and obstruction cohomology,
    Church--Rosser confluence, game-theoretic iterated admissibility
    (Bicchieri--Schulte), HoTT transport, explicit substitution calculi.
    These are established mathematical frameworks.
    The generalized reduction operator $\Rcal$ (Definition~\ref{def:R})
    lives at this layer.

  \item[\textbf{Layer~IV: Philosophical synthesis.}]
    Hofstadterian analogy and chunking,
    Spherepop membrane collapse as geometric operational semantics,
    RSVP thermodynamic cognition,
    observer admissibility as meta-constraint on coherent existence
    (Modgil--Patil, Zenteno, Paton).
    These are interpretive proposals organizing the lower layers.
    The admissibility principle as stated in this monograph lives
    primarily here.

  \item[\textbf{Layer~V: Research program.}]
    Define generalized reduction operators with explicit soundness
    proofs; develop admissibility metrics for each domain;
    investigate category-theoretic transport structures (natural
    transformations, fibered categories) capable of relating
    admissibility conditions across layers; formalize the
    generalized confluence schema; calibrate RSVP parameters
    from first principles or experiment; validate Marine/Phoenix 8a
    against established audio benchmarks.
\end{description}

The stratification protects against the two failure modes identified
in the introduction.
It prevents premature literalism by clearly separating Layer~I
(empirically established) from Layers~III--IV (formal/philosophical).
It prevents vacuous analogy by requiring, at Layer~III, explicit
mathematical content for each instantiation of $\Rcal$.

\begin{table}[h!]
\centering
\caption{Extended concept atlas with epistemic layer.}
\label{tab:extended-atlas}
\begin{tabular}{llll}
\toprule
\textbf{Domain} & \textbf{Admissibility detector} & \textbf{Failure mode}
  & \textbf{Layer} \\
\midrule
GR / vacuum      & Kretschmann scalar $K$ & Curvature singularity
  & I / IV \\
$\lambda$-calculus & Type preservation $\Gamma\vdash M:T$ & Non-termination
  & III \\
Sheaf semantics   & $\check{H}^1 = 0$ & Gluing obstruction
  & III \\
RSVP              & $\dot{S} \leq \dot{\Sigma}_{\mathrm{crit}}$ & Supercritical runaway
  & IV \\
Turbulence        & Maximum dissipation rate & Sub-optimal $\alpha_b$
  & I \\
Spherepop         & Historical membrane coherence & Infinite bubbling
  & IV \\
MEM$|$8           & Resonant retrieval boundedness & Attractor fragmentation
  & II \\
Hofstadter        & Analogical coherence preservation & Recursive reminding
  & IV \\
Game theory       & Iterated weak admissibility (IWD) & Dominated strategy
  & III \\
Observer (GR)     & $g_{\mu\nu}V^\mu V^\nu > 0$ & Trans-Planckian divergence
  & IV \\
Marine / Phoenix  & Harmonic-lock stability & Lock failure / noise floor
  & II \\
\bottomrule
\end{tabular}
\end{table}

\subsection{The invariant that unifies all layers}

Across all eleven domains in Table~\ref{tab:extended-atlas}, the
same structural claim recurs:

\begin{quote}
  \textbf{Stable systems preserve admissibility under local
  transformation.}
\end{quote}

In $\lambda$-calculus, $\beta$-reduction preserves typing:
\[
  \Gamma \vdash (\lambda x.M)N : T
  \;\Rightarrow\;
  \Gamma \vdash M[x:=N] : T.
\]
In Spherepop, membrane collapse preserves historical topology.
In Hofstadter, analogy preserves conceptual coherence under
contextual rebinding.
In RSVP, entropy gradients preserve viable future trajectories.
In sheaf semantics, global sections exist only when
$\check{H}^1(\mathcal{G}_\infty) = 0$.
In quantum optics, interference preserves phase admissibility.
In turbulence, maximum dissipation preserves the physically correct
mixing trajectory.
In game theory, iterated elimination preserves
representation-independence of the solution.
In relativistic observer theory, local Lorentz symmetry is preserved
exactly within the admissible domain.

The admissibility principle is the claim that this invariant-preservation
structure is not accidental but reflects a common operational grammar.

% ================================================================
\section{Admissibility and Constraint-Limited Observerhood:
  RSVP, CLIO, TARTAN, and Yarncrawler}
\label{sec:rsvp-synthesis}
% ================================================================

\begin{speculative}
  This section integrates the admissibility program with several
  Flyxion frameworks (RSVP, CLIO, TARTAN, Yarncrawler) that are
  under active theoretical development.
  The connections drawn here are interpretive proposals, not
  established results.
\end{speculative}

\subsection{Observerhood as emergent admissibility}

The admissibility framework developed across the preceding sections
aligns naturally with the RSVP program, the CLIO
(\emph{Constraint-Leveraged Inference Operator}) formalism, and the
projection-based interpretation of cognition and physics developed
throughout this work.

In RSVP, physical systems are not isolated objects moving through a
passive geometric background.
All observable structure emerges from constraint evolution within the
scalar--vector entropy field $(\Phi, \mathbf{v}, S)$, where $\Phi$
denotes scalar coherence density, $\mathbf{v}$ directed constraint
flow, and $S$ admissible configurational entropy.
Observerhood itself becomes a dynamically stabilized trajectory through
constraint space.
An observer is therefore not fundamental but emergent from persistent
low-entropy coherence structures capable of recursively maintaining
informational continuity.

The admissibility principle acquires a reformulation in RSVP language:
\begin{quote}
  An observer is physically admissible if and only if its trajectory
  remains recursively coherent under the local and global constraint
  geometry of the plenum.
\end{quote}

The relativistic admissibility condition $g_{\mu\nu}V^\mu V^\nu > 0$
becomes only the lowest-order projection of a more general coherence
requirement.
The metric itself emerges as a coarse-grained geometric encoding of
admissible informational trajectories within the plenum.

\subsection{CLIO: cognition as constrained projection}

In CLIO, cognition and measurement are sparse projection operators
acting on inaccessible high-dimensional state spaces.
Let
\[
  \pi : X \to M
\]
denote the projection from a high-dimensional ontic trajectory manifold
$X$ into an operational measurement manifold $M$.
The observer navigates compressed projections preserving only
dynamically relevant constraint relations.
Admissibility requires not merely geometric timelikeness but
\emph{projectional stability}:
\[
  \pi(R(x)) = \pi(x)
\]
for admissible reduction operators $R$ preserving coherent inferential
structure.

This clarifies why Planck-scale divergences arise operationally rather
than ontologically.
Attempts to probe arbitrarily small scales force the observer to
collapse the projection manifold faster than recursive coherence can
be maintained.
The failure is not a breakdown of spacetime itself but a breakdown of
observer-stabilized projection consistency.

Entropy in this reading is not merely disorder but the logarithmic
volume of dynamically accessible future trajectories:
\[
  S \sim \log |\mathcal{A}|,
\]
where $\mathcal{A}$ is the admissible future trajectory set.
Infinite boosts, singular localization, and trans-Planckian measurement
attempts collapse $|\mathcal{A}|$ toward instability.

\subsection{TARTAN and Yarncrawler}

TARTAN (\emph{Trajectory-Aware Recursive Tiling with Annotated Noise})
treats physical and cognitive systems as recursively tiled trajectory
bundles carrying both dynamic evolution and semantic metadata.
Observerhood corresponds to maintaining stable trajectory alignment
across recursively compressed scales.
Admissibility becomes equivalent to preserving trajectory continuity
under multiscale coarse-graining.

Yarncrawler extends this to civilization-scale systems.
Roads, infrastructures, memory systems, and institutions are
trajectory stabilizers preserving admissible movement through physical
and semantic space.
A city ceases functioning not when matter disappears but when transport
pathways lose coherent accessibility.
Maintenance itself is a form of admissibility preservation:
repair systems, garbage systems, and archival systems prevent collapse
into inadmissible trajectory fragmentation.

\subsection{The philosophical inversion}

The synthesis yields a central philosophical inversion of the RSVP
program:
\begin{quote}
  Physical law does not merely govern what exists;
  it governs what can remain coherently observable.
\end{quote}

Admissibility is not an additional force or symmetry.
It is a meta-constraint on coherent existence itself.

More precisely: the fundamental question of physics changes.
Instead of asking only what mathematically exists, we begin asking
what can remain coherently observable.
The universe becomes not merely a collection of possible states but a
dynamic filter selecting which trajectories can persist without
collapsing their own conditions of observation.

This reformulation naturally integrates:
\begin{itemize}
  \item relativistic geometry (admissibility as timelikeness and
        curvature boundedness);
  \item quantum operational limits (Planck scale as projectional
        collapse horizon);
  \item recursive cognition (thought as admissibility-maintenance
        under analogical pressure);
  \item informational thermodynamics (entropy as trajectory
        accessibility volume);
  \item social systems (institutions as admissibility-preserving
        trajectory stabilizers);
  \item semantic systems (meaning as globally glueable local
        consistency).
\end{itemize}

% ================================================================
\section{Discussion and Open Problems}
\label{sec:conclusion}
% ================================================================

\subsection{What has been established}

The following claims in this paper rest on established results:
\begin{itemize}
  \item The $\lambda$-calculus formalism, $\beta$-reduction,
        Church--Rosser confluence, and Curry--Howard are standard
        (Church 1936; Sørensen--Urzyczyn 2006).
  \item The quantum dwell-time results (anomalous weak values,
        group-delay formulas) follow APL Quantum (2025).
  \item The DC-OPA performance figures follow the Nature Photonics
        report on the advanced BiBO/MgO:LN system.
  \item Hofstadter's theory of analogy and chunking is presented
        faithfully from published work (1979, 1995, 2013).
  \item The sheaf-theoretic gluing and obstruction framework is
        standard category theory (Kashiwara--Schapira 1994).
  \item The game-theoretic admissibility results (iterated weak
        dominance, representation-independence) follow
        Bicchieri--Schulte (1996).
  \item The turbulence validation results (maximum dissipation rate,
        $\alpha_b \approx 0.06$ matching RT experiment) follow
        Glimm et al.\ (preprint LA-UR-19-20285).
  \item The Kretschmann scalar as minimal non-degenerate curvature
        invariant is standard differential geometry.
  \item The multimodal lineage tracing methodology, computational
        frameworks (LineageOT, moslin, LineageVAE, PORCELAN), and
        characterization of progressive potency restriction follow
        Wang, He, and Hu (2026) in \textit{Nature Reviews Genetics}.
        The interpretation of lineage tracing through the admissibility
        schema is the authors' proposal.
  \item The underconstraint argument follows Yudkowsky (2008);
        the endogenous growth and Moore's Law analyses are as
        described in that essay.
\end{itemize}

\subsection{What is proposed but not yet established}

\begin{itemize}
  \item The Spherepop--$\lambda$-calculus formal equivalence requires a
        complete syntactic specification of Spherepop and a soundness
        proof for Table~\ref{tab:lambda-spherepop}.
  \item The RSVP model parameters ($\lambda_c$, $\dot{\Sigma}_{\mathrm{crit}}$)
        require empirical calibration or derivation from first principles.
  \item The Marine/Phoenix 8a architecture requires independent benchmark
        evaluation, ablation studies, and perceptual testing.
  \item The generalized confluence schema (Schema~1) is a conjecture;
        it holds in specific cases but has not been proved in general.
  \item The Hofstadter--Spherepop cognitive mapping is interpretive;
        it provides a framework for future empirical and formal work,
        not a proof.
  \item The observer-admissibility framework of Modgil--Patil (2026)
        requires covariant reconstruction of $\gamma_g$ using tetrads
        or observer congruences; the global-bound formula is
        currently phenomenological.
  \item The Zenteno (2026) penalty-functional approach requires
        derivation of the penalty from an action principle.
\end{itemize}

\subsection{Priority of the research program}

Layer~V of the five-layer architecture (Section~\ref{sec:fivelayers})
constitutes the forward research agenda.
In order of logical priority:
\begin{enumerate}
  \item Formalize Spherepop as an explicit substitution calculus and
        prove the correspondence with $\lambda\sigma$ or $\lambda\upsilon$.
  \item Prove or refute the generalized confluence schema
        (Schema~1) for the RSVP relaxation operator under explicit
        stability assumptions.
  \item Develop a covariant tetrad-based formulation of the
        generalized relativistic factor $\gamma_g$ to put the
        observer-admissibility framework on rigorous footing.
  \item Investigate whether the Kretschmann-scalar admissibility
        detector admits categorical generalization:
        what is the analogue of $K$ for semantic, cognitive, and
        thermodynamic admissibility?
  \item Apply the Bicchieri--Schulte recursion to multi-agent
        semantic systems; formalize whether LLM hallucination rates
        decrease under iterated admissibility elimination over
        agents' epistemic states.
  \item Develop independent benchmarks for the Marine
        temporal-interval approach against STFT baselines on
        standardized audio corpora.
  \item Formalize the sheaf-hallucination proposition and investigate
        testable predictions about failure modes in language models.
  \item Explore whether Hofstadterian cognitive admissibility can be
        operationalized in a computational architecture (e.g.,
        active inference or predictive coding frameworks).
\end{enumerate}

\subsection{Admissibility as research grammar rather than universal law}

The reviewer of an early draft of this monograph offered the following
observation: the paper's closing remark --- that the admissibility
principle is ``a research grammar'' rather than a physical law or
mathematical theorem --- is the paper's most accurate statement of
its own contribution.
That observation is correct, and the present work embraces rather
than retreats from it.

The comparison to category theory is apt and worth dwelling on.
Category theory did not claim that all mathematics was literally the
same object, or that groups and topological spaces were formally
equivalent.
It supplied a \emph{transport language} for structural relationships
across domains: functors, natural transformations, adjoints, limits.
This language proved extraordinarily productive not because it proved
old theorems but because it made new questions visible.

The admissibility principle aims to occupy the same role.
Its contribution is not a theorem but a vocabulary:
admissibility region, perturbation, collapse operator, invariant,
normal form, confluence, gluing obstruction.
These terms, applied with the precision demanded by
Definition~\ref{def:R} and the necessary conditions in the preceding
subsection, allow researchers to ask coherent questions about
constraint-preserving collapse in domains where those questions had
not previously been posed in that form.

The most defensible version of the central thesis is therefore:

\begin{quote}
  Many systems --- formal, physical, cognitive, social --- exhibit
  recurring structural motifs involving constrained trajectory
  reduction, invariant preservation, and local compatibility
  propagation.
  These motifs are sufficiently precise to motivate a shared research
  grammar.
  Identifying the grammar is a contribution.
  Proving its instances are formally equivalent is the research
  program that follows from it.
\end{quote}

% ================================================================
\section{Underconstrained Abstractions and the Limits of Historical Fit}
\label{sec:underconstraint}
% ================================================================

\subsection{The problem of historical underconstraint}

One of the deepest epistemic problems facing any cross-domain synthesis is that
historical data radically underdetermines abstraction.
A single body of observations may admit many structurally compatible descriptions, and
the mere existence of a retrospective fit does not establish that an abstraction
captures the operative mechanism responsible for the observed behavior.
A curve can be interpolated through finitely many points by polynomials of arbitrarily
different character; the interpolation tells us nothing about which polynomial governs
the generative process.

This problem was articulated with particular sharpness by Eliezer Yudkowsky in a 2008
essay titled ``Underconstrained Abstractions,'' written as part of an extended exchange
with Robin Hanson about the epistemic status of abstractions used in technological
forecasting.
Yudkowsky's central observation is worth quoting in its original form: ``speaking of an
abstraction being `verified' by previous history is a tricky thing.
There is this little problem of underconstraint --- of there being more than one
possible abstraction that the data `verifies'.''
He continues: ``The further away you get from highly regular things like atoms, and the
closer you get to surface phenomena that are the final products of many moving parts,
the more history underconstrains the abstractions that you use.
This is part of what makes futurism difficult.
If there were obviously only one story that fit the data, who would bother to use
anything else?''

The problem is structural rather than merely practical.
In domains close to fundamental physics, the regularity of the underlying processes
tightly constrains the space of compatible abstractions.
Atomic physics admits few retrospective stories because the entities involved obey
highly symmetric and repeatable laws that sharply penalize deviant descriptions.
But as one ascends toward emergent, historically layered systems --- cognition,
economics, institutional dynamics, semantic systems, civilization-scale technological
development --- the surface phenomena are the final product of many interacting layers,
each of which introduces additional degrees of freedom that the historical record does
not constrain.
In such domains, retrospective fit is necessary but very far from sufficient.
An abstraction that merely redescribes historical data without imposing constraints on
what patterns it excludes is, in Yudkowsky's sense, underconstrained.

\subsection{The endogenous growth illustration}

Yudkowsky's primary illustration of the underconstraint problem involves the
``endogenous growth'' literature in economics, specifically theories in which knowledge
accumulation is modeled as a combinatorial process.
The seminal formulation treats ideas as being generated by combining other ideas $N$ at
a time, so that a stock of $N$ available ideas produces on the order of $N!/(k!(N-k)!)$
potential new ideas for any fixed combination size $k$.
From this combinatorial abundance, it follows that the number of available ideas will
vastly exceed the number of people available to exploit them, and the rate of knowledge
production will therefore be limited by the supply of researchers rather than by the
supply of ideas.

Yudkowsky makes a precise and important observation about this structure.
If the only proposition that actually enters the predictive machinery of the theory is
that there exist more potential ideas than people to pursue them, then this is the only
proposition that the theory's fit to historical data can even partially verify.
The elaborate combinatorial apparatus — the $N$-choose-$k$ framing, the specific
functional form of the idea-generation process — is not doing any predictive work
beyond what would be accomplished by the simpler statement: ``I assume there are more
ideas than people to exploit them.''
Any hypothesis with ``more ideas than people to exploit them'' in its structure will
fit the same historical data equally well.
The additional mathematical specificity is not constrained by the data because it is
not used in the forward application; it provides a causal-sounding story without
providing causal leverage.

This is the core epistemic failure Yudkowsky identifies: an abstraction can achieve
the appearance of precision and the appearance of historical verification while
remaining genuinely underconstrained, because the constrained parts of the abstraction
are not the parts doing the predictive work.
The decorative mathematics acquires false credibility by association with the parts of
the theory that are actually tested.

\subsection{Moore's Law and multiple compatible abstractions}

The Moore's Law case is the most concrete illustration in Yudkowsky's essay, and it
bears careful examination because it directly concerns the type of technological
forecasting that has the highest stakes for arguments about artificial intelligence.

The historical trajectory of exponential growth in transistor density is compatible
with at least four distinct generating abstractions, all fitting the same observed
data.
One abstraction treats transistor count as a function of elapsed calendar time, an
empirical regularity with no identified causal mechanism beyond the observation itself.
A second treats it as a consequence of exponential growth in capital investment in
semiconductor manufacturing and research and development, predicting continuation as
long as that investment continues.
A third treats it as a function of optimization pressure applied by a community of
human researchers operating at a fixed cognitive rate, predicting continuation as long
as the research community remains adequately populated and funded.
A fourth treats it as involving some degree of recursive amplification, in which
improved tools make further improvement easier, predicting a non-linear relationship
between the trajectory of improvement and the substrate on which the researchers
themselves operate.

Yudkowsky raises the counterfactual that makes these abstractions empirically
distinguishable in principle while remaining indistinguishable within the historical
record: ``Would the further progress of Moore's Law have been different from that in
our own world, relative to sidereal time, if, for these last fifty years, the
researchers themselves had been running on the latest generation of computer chip at
any given point?''
This counterfactual is not answerable from historical data, because the researchers
were not in fact running on progressively faster hardware during the period when the
historical data was generated.
But the four competing abstractions give sharply different answers to it, and those
different answers imply radically different predictions about what would happen if
the researchers' cognitive substrate were to change.

The underconstrained-abstractions problem is therefore not merely that multiple stories
fit the data.
It is that the stories diverge precisely where they matter most for forward application,
and the historical data provides no basis for choosing between them in that region.

\subsection{Moravec and the forward-application problem}

Yudkowsky's essay also references Moravec's 1988 work as an illustration of a
complementary failure: the problem of assuming a unique forward application of an
abstraction that has, in fact, multiple possible extensions.
Moravec estimated the computing power required for human-level artificial intelligence
by multiplying the apparent number of operations performed by the human brain per second
by an appropriate efficiency factor, arriving at a figure he then extrapolated forward
using Moore's Law to estimate the arrival date.
Yudkowsky notes that Moravec ``spent a lot of time talking about how much `computing
power' the human brain seems to use --- but much less time talking about whether an AI
would use the same amount of computing power, or whether using Moore's Law to
extrapolate the first supercomputer of this size is the right way to time the arrival
of AI.''

The problem is not that Moravec's historical estimates of brain computational
equivalence were wrong, but that the forward application contained a hidden assumption:
that an artificial system performing at human level would require approximately the
same number of operations per second as the biological system being emulated, and that
the arrival of sufficient hardware therefore determines the arrival of the capability.
This assumption is not constrained by the historical data used to establish the
computational-equivalence estimate.
It is an additional hypothesis that enters silently at the point of forward application,
and it happens to be the hypothesis doing all the predictive work in the arrival-date
calculation.
The estimate is therefore not a prediction so much as a restatement of the hidden
assumption, dressed in the credibility borrowed from the carefully validated
computational-equivalence estimate.

This is what Yudkowsky calls the problem of ``a lot of slack with regards to exactly
what will happen with respect to that abstraction, going forward'': after constructing
a historical story that appears to verify an abstraction, the forward application
typically requires additional choices about how to apply the abstraction, and those
choices are not themselves constrained by the historical verification.
The appearance of a single determinate prediction conceals an underspecified space of
possible applications.

\subsection{The admissibility response to underconstraint}

The admissibility framework attempts to respond to this problem not by eliminating
abstraction but by constraining which kinds of abstractions count as legitimate
transport structures.

An admissible abstraction must do more than fit historical data.
It must preserve identifiable invariants under perturbation, specify failure modes
precisely enough that they could in principle be observed, restrict future trajectories
in ways that generate differential predictions across competing causal stories, and
produce exclusions rather than merely redescriptions.
An abstraction that imposes no exclusions is vacuous; it is capable of absorbing any
observation after the fact by adjusting its parameters or reinterpreting its terms.
The formal apparatus of Definition~\ref{def:R} and the augmented strong admissibility
criterion of Appendix~\ref{app:nonvacuity} are therefore not merely mathematical
decoration.
They are the mechanisms by which the framework resists the slide into explanatory
vacuity.

The connection to Yudkowsky's diagnosis is direct.
An underconstrained abstraction in his sense is precisely an abstraction that fails to
specify which of its components are doing the predictive work and which are providing
a causal-sounding story without causal leverage.
The admissibility framework operationalizes this distinction by requiring that the
collapse operator, the compatibility condition, and the invariant be made explicit
and distinct.
An abstraction in which the ``invariant'' is never operationalized, in which the
``compatibility condition'' is satisfied by any input, or in which the ``collapse
operator'' is never applied to generate a prediction that could fail, is an abstraction
that fails the admissibility criterion.
It may fit the history, but it is not doing so through a mechanism that constrains the
future.

\subsection{The danger of unconstrained universality for the present framework}

This diagnosis applies reflexively to the admissibility principle itself.
If the admissibility framework were to become underconstrained in Yudkowsky's sense,
it would fit all observed historical instances of ``constraint-preserving
transformation'' while providing no leverage on which future processes will or will
not exhibit admissibility structure.
The framework would then be exactly the kind of abstraction it is designed to
discipline.

The non-vacuity criterion of Appendix~\ref{app:nonvacuity} and the three augmented
conditions of strong admissibility are the framework's attempt to avoid this failure.
Requiring that the admissibility predicate recursively restructure the future trajectory
space, that the invariant be transportable and usable by downstream processes, and that
the system exhibit genuine path-sensitivity are requirements specifically designed to
ensure that the admissibility framework is not merely a retrospective vocabulary.
A thermostat satisfies a liberal reading of Definition~\ref{def:R} but fails strong
admissibility precisely because no discriminatory work is being done: the thermostat's
future behavior is fully determined by its instantaneous state, there is no invariant
being transported downstream, and there is no path-sensitivity.
Excluding the thermostat is not an embarrassment; it is the framework demonstrating
that it imposes genuine constraints.

The same logic applies to the cross-domain cases treated throughout this monograph.
The sheaf-theoretic hallucination model, the Glimm turbulence admissibility criterion,
the Bicchieri--Schulte recursive epistemic admissibility, and the Kretschmann scalar
curvature detector are all cases where the admissibility structure is specified with
enough precision to be wrong.
The RSVP framework and Spherepop are cases where the specification is not yet
sufficiently precise, and the monograph identifies them explicitly as research programs
rather than established results.
This stratification is the monograph's attempt to practice what it preaches: not all
cross-domain analogies are equally constrained, and the framework's own epistemic
standards require acknowledging the difference.

\subsection{Forecasting as constraint-topology tracking}

The future is difficult to predict not merely because it has not yet occurred but
because the admissibility manifold may itself change.
Historical regularities emerge from fixed constraints, contingent optimization
pressures, hidden mesoscale structures, or temporary equilibrium conditions, and when
any of these underlying structures shifts, an abstraction previously fitted to the
historical trajectory may cease to apply not gradually but catastrophically.

This is why extrapolation across phase transitions in the admissibility manifold is
unreliable even for abstractions with excellent retrospective fit.
A system near a bifurcation point in its constraint topology generates historical data
that looks regular and predictable, because the constraint structure has been stable
throughout the observed period.
But an abstraction fitted to that period carries no information about the system's
behavior after the bifurcation.
The admissibility principle therefore treats forecasting as a problem of tracking
constraint evolution rather than extending historical curves.
Yudkowsky's question about Moore's Law in the counterfactual world where researchers
run on their own products is precisely the question of whether the constraint topology
would have changed, and if so, how the abstraction would need to be modified to remain
valid under the new topology.
Answering it requires understanding what generates the observed regularities, not
merely that they have been observed.

\subsection{Connection to sheaf-theoretic obstruction}

The problem of underconstrained abstraction is closely related to the sheaf-theoretic
notion of gluing failure developed in Section~\ref{sec:unification} and
Appendix~\ref{app:sheaf}.

An abstraction may remain locally compatible with historical data while failing to
extend coherently across perturbative or counterfactual domains.
The historical record constitutes a collection of local sections $\{s_i\}$ over an
open cover of the observation domain, and the abstraction proposes that these sections
glue into a global section extending beyond the observed region.
The existence of pairwise-compatible local sections does not guarantee the existence
of a global section: the obstruction to extension is measured by the first
\v{C}ech cohomology class $[\omega] \in \check{H}^1(X, \mathcal{F})$.

When $[\omega] \neq 0$, the abstraction glues coherently over the observed history
but fails globally once the constraint topology changes.
This is the sheaf-theoretic reading of Yudkowsky's underconstraint problem: the
abstraction satisfies the cocycle condition on observed overlaps but cannot be extended
to a globally coherent section of the semantic manifold.
The forward-application failures Yudkowsky identifies --- Moravec's computational
extrapolation, endogenous growth models applied to recursive self-improvement,
Moore's Law applied across a change in the cognitive substrate of researchers --- are
all cases where local historical compatibility does not imply global extension.
The obstruction is invisible within the range of observations used to construct the
abstraction and becomes apparent only in the counterfactual or future domain where the
abstraction is actually needed.

The admissibility framework offers a partial response.
By requiring that an admissible abstraction specify its failure modes explicitly, it
forces the question of where the global section might fail to extend.
An abstraction that specifies the conditions under which its invariant ceases to be
preserved, the boundary of its admissibility region, and the perturbations that fall
outside its compatibility condition is an abstraction whose obstruction classes are at
least partially identified before the failure occurs.
This does not eliminate the problem of underdetermination, but it converts a silent
failure mode into a visible research question: whether the abstraction extends globally
becomes a question with specific mathematical content rather than an implicit assumption
whose failure is discovered only after the extrapolation has been committed to.

\subsection{Closing remark}

The admissibility principle is not a physical law, a mathematical
theorem, or an empirical discovery.
It is a \emph{research grammar}: a way of asking the same question
across very different systems.
That question is: \emph{what makes a perturbation locally receivable,
and what becomes stable once it has been received?}

When that question can be answered precisely for a given domain,
the admissibility schema becomes a theorem.
When it can be answered approximately, it becomes a model.
When the answer remains unclear, it remains heuristic.

The purpose of this monograph is to make the question precise enough
that the distinction between these three cases can be drawn.

% ================================================================
% Bibliography (manual; replace with BibTeX entries as desired)
% ================================================================
\newpage
\section*{References}
\addcontentsline{toc}{section}{References}

\begin{itemize}[leftmargin=1.8em, itemsep=0.4em]

  \item Church, A. (1936). An unsolvable problem of elementary number
        theory. \textit{American Journal of Mathematics}, 58(2), 345--363.

  \item Curry, H.~B., \& Feys, R. (1958). \textit{Combinatory Logic},
        Vol.~I. North-Holland.

  \item Girard, J.-Y., Taylor, P., \& Lafont, Y. (1989).
        \textit{Proofs and Types}. Cambridge University Press.

  \item Hofstadter, D.~R. (1979). \textit{Gödel, Escher, Bach:
        An Eternal Golden Braid}. Basic Books.

  \item Hofstadter, D.~R., \& Mitchell, M. (1995).
        \textit{Fluid Concepts and Creative Analogies}.
        Basic Books.

  \item Hofstadter, D.~R., \& Sander, E. (2013).
        \textit{Surfaces and Essences: Analogy as the Fuel and Fire
        of Thinking}. Basic Books.

  \item Howard, W.~A. (1980). The formulae-as-types notion of
        construction. In J.~P. Seldin \& J.~R. Hindley (Eds.),
        \textit{To H.~B. Curry: Essays on Combinatory Logic, Lambda
        Calculus, and Formalism} (pp. 479--490). Academic Press.

  \item Sørensen, M.~H., \& Urzyczyn, P. (2006).
        \textit{Lectures on the Curry--Howard Isomorphism}.
        Elsevier (Studies in Logic and the Foundations of Mathematics,
        Vol. 149).

  \item [APL Quantum] Author(s). (2025). Atomic excitation times
        for transmitted and scattered photons.
        \textit{APL Quantum}, \textbf{2}, 036108.
        [Full author list to be inserted.]

  \item [Nature Photonics DC-OPA] Author(s). (2024--25).
        Advanced dual-chirped optical parametric amplification
        producing single-cycle MIR pulses.
        \textit{Nature Photonics}.
        [Full citation to be inserted.]

  \item Kashiwara, M., \& Schapira, P. (1994).
        \textit{Sheaves on Manifolds}. Springer.

  \item de Bruijn, N.~G. (1972). Lambda calculus notation with
        nameless dummies. \textit{Indagationes Mathematicae},
        34(5), 381--392.

  \item Abadi, M., Cardelli, L., Curien, P.-L., \& Lévy, J.-J.
        (1991). Explicit substitutions. \textit{Journal of Functional
        Programming}, 1(4), 375--416.

  \item Bicchieri, C., \& Schulte, O. (1996). Common reasoning about
        admissibility. \textit{Erkenntnis} (revised 1996, forthcoming).

  \item Evangelista, R. (preprint). Admissibility and the minimal
        curvature bound in covariant theories. [Author affiliation not
        given; full citation to be completed.]

  \item Glimm, J., Cheng, B., Sharp, D.~H., \& Kaman, T.
        A crisis for the verification and validation of turbulence
        simulations. Los Alamos National Laboratory Preprint
        LA-UR-19-20285. [Published version citation to be inserted.]

  \item Johansson, I. (year). Construction as reduction.
        [Full citation to be inserted --- referenced for the analogy-as-reduction
        framing in linguistics.]

  \item Hawking, S.~W. (1975). Particle creation by black holes.
        \textit{Communications in Mathematical Physics}, 43, 199--220.

  \item Bekenstein, J.~D. (1973). Black holes and entropy.
        \textit{Physical Review D}, 7, 2333--2346.

  \item Bousso, R. (2002). The holographic principle.
        \textit{Reviews of Modern Physics}, 74, 825--874.

  \item Jacobson, T. (1995). Thermodynamics of spacetime: the Einstein
        equation of state.
        \textit{Physical Review Letters}, 75, 1260--1263.

  \item de Lellis, C., \& Székelyhidi, L. (2009).
        The Euler equations as a differential inclusion.
        \textit{Annals of Mathematics}, 170, 1417--1436.
        [Cited for wild Euler-equation solutions establishing
        nonuniqueness; background for the Glimm et al.\ turbulence
        admissibility argument.]

  \item Modgil, M.~S., \& Patil, D.~D. (2026). On the invariance of
        Planck scale under special relativity and the geometry of
        quantum space. Cosmos Research Labs. [Preprint; February 9, 2026.]

  \item Paton, A.~J. (2026). Structural admissibility as a
        pre-explanatory condition in scientific reasoning.
        The Paton System.

  \item Wang, K., He, X., \& Hu, Z. (2026).
        Computational approaches for multimodal lineage tracing.
        \textit{Nature Reviews Genetics}.
        \doi{10.1038/s41576-026-00969-9}
        [Primary source for the biological admissibility instantiation
        in Section~\ref{sec:lineage}.]

  \item Yudkowsky, E. (2008). Underconstrained Abstractions.
        \textit{LessWrong} (December 4, 2008).
        \url{https://www.lesswrong.com/posts/underconstrained-abstractions}
        [Cited for the problem of historical underconstraint in technological
        forecasting, the endogenous growth illustration, and the Moore's Law
        multiple-abstraction analysis.]

  \item Moravec, H. (1988). \textit{Mind Children: The Future of Robot and Human
        Intelligence}. Harvard University Press.
        [Cited for the computational-equivalence extrapolation and forward-application
        problem discussed in Section~\ref{sec:underconstraint}.]

  \item Zenteno, C. (2026). Admissibility under constraint: dynamical
        criterion beyond existence and consistency.
        Zoa Industries --- Compressed Consciousness Laboratory.
        Zenodo. \url{https://doi.org/10.5281/zenodo.18370414}


\end{itemize}

% ================================================================
\appendix
% ================================================================

% ================================================================
\section{Category-Theoretic Formulation of Admissibility}
\label{app:category}
% ================================================================

This appendix develops a category-theoretic formulation of the generalized reduction
operator introduced in Section~\ref{sec:unification}. The goal is not to assert a
completed categorical foundation for the admissibility principle, but to identify
the minimal mathematical structures such a foundation would require, and to demonstrate
that the framework is at least as rich as established categorical machinery in several
limiting cases.

\subsection{Admissibility categories}

Let $\mathcal{C}$ be a category whose objects are admissibility regions $\Bcal$ and
whose morphisms represent admissibility-preserving transformations between them.
We equip $\mathcal{C}$ with three additional pieces of structure.

The first is a partial reduction endofunctor
\[
  \mathfrak{R} : \mathcal{C} \rightharpoonup \mathcal{C},
\]
defined only on those objects $\Bcal_i$ for which there exists an admissible
perturbation $\Delta_i$ compatible with the constraint structure $\partial(\Bcal_i)$.
Partiality here is essential: the functor is not globally defined, and the domain
of $\mathfrak{R}$ is determined by the admissibility predicate rather than by
universal coverage.

The second piece of structure is a family of admissibility predicates
$\mathsf{Adm}_{A,B}(f)$ on morphisms $f : A \to B$, specifying which morphisms
participate in coherent collapse sequences.
A morphism $f$ is admissible if the image $\mathfrak{R}(A)$ is reachable from $A$
via $f$ without violating the invariant condition.

The third is an invariant-preservation natural transformation
\[
  \eta : \iota \Rightarrow \iota \circ \mathfrak{R},
\]
where $\iota : \mathcal{C} \to \mathcal{I}$ is a functor into an invariant category
$\mathcal{I}$ tracking the transportable structure that collapse must preserve.
The naturality of $\eta$ encodes the requirement that invariant preservation is
uniform across all admissible reductions: no reduction may preserve the invariant
selectively for some morphisms while violating it for others.

\begin{definition}[Admissibility category]
  An \emph{admissibility category} is a tuple
  $(\mathcal{C}, \mathfrak{R}, \mathsf{Adm}, \iota, \eta)$ in which $\mathcal{C}$ is
  a category, $\mathfrak{R}$ is a partial endofunctor, $\mathsf{Adm}$ is a family of
  admissibility predicates on morphisms, $\iota$ is an invariant functor, and $\eta$
  is a natural transformation $\eta_A : \iota(A) \to \iota(\mathfrak{R}(A))$
  for every $A$ in the domain of $\mathfrak{R}$, with each $\eta_A$ an isomorphism
  in $\mathcal{I}$.
\end{definition}

The isomorphism condition on $\eta_A$ captures the idea that admissible reduction does
not merely approximate the invariant but preserves it exactly. In typed
$\lambda$-calculus, this corresponds to the Subject Reduction theorem:
$\beta$-reduction preserves typing judgements, so the type is the invariant and the
preservation is exact under the isomorphism $\eta_A : T_A \xrightarrow{\sim} T_{A'}$.
In sheaf semantics it corresponds to the requirement that gluing maps are isomorphisms
on overlapping sections.

\subsection{Ordinary $\lambda$-calculus as a special case}

The untyped $\lambda$-calculus corresponds to the special case in which the
admissibility predicate is trivially satisfied after $\alpha$-conversion: every redex
is admissible, $\mathfrak{R}$ is defined everywhere on reducts, and the invariant
category $\mathcal{I}$ is trivial, reflecting the fact that untyped reduction carries
no non-trivial invariant beyond the reduction relation itself.

In the simply-typed setting, $\mathcal{I}$ becomes the category of simple types and
type-preserving substitutions, and Subject Reduction provides $\eta$. In System~F,
$\mathcal{I}$ becomes the category of polymorphic types. In dependent type theory,
$\mathcal{I}$ becomes a category whose objects are dependent types and whose morphisms
are definitional equalities. The admissibility framework thus interpolates cleanly
across the standard hierarchy of type theories, interpreting each as an admissibility
category with progressively richer invariant structure.

Spherepop generalizes this by making admissibility depend on historical topology.
The object $A = (\mathcal{I}(A), \partial(A), H(A))$ in an admissibility category for
Spherepop carries a third component, the history $H(A)$, which enters the admissibility
predicate: $\mathsf{Adm}_{A,B}(f)$ holds only if $f$ is consistent with $H(A)$ as a
causal continuation. This is the categorical expression of condition~(iii) in
Definition~\ref{def:R}.

\subsection{Confluence as a weak diamond condition}

The Church--Rosser property admits a clean categorical reformulation. In the
admissibility category, confluence becomes the following.

\begin{definition}[Weak admissibility diamond]
  An admissibility category $(\mathcal{C}, \mathfrak{R}, \mathsf{Adm}, \iota, \eta)$
  satisfies the \emph{weak admissibility diamond property} if for every pair of
  admissible morphisms $f : A \to B$ and $g : A \to C$ in $\mathcal{C}$, there
  exists an object $D$ together with admissible morphisms $f' : B \to D$ and
  $g' : C \to D$ such that $f' \circ f = g' \circ g$ holds up to the equivalence
  relation generated by the admissibility predicates.
\end{definition}

This definition generalizes Church--Rosser confluence, sheaf gluing conditions,
coherence diagrams in monoidal categories, and path reconciliation in homotopy theory.
In each case the weak admissibility diamond asserts that locally incompatible collapse
paths can be joined into a common successor. These are not the same theorem: each
domain proves the diamond by entirely different arguments. The admissibility category
framework supplies a common \emph{language} in which the analogous structural
condition in each domain is stated in the same syntactic form.

\subsection{Functorial transport between admissibility categories}

A central aspiration of the monograph is that admissibility structure can be
transported between domains. In categorical terms this requires functors
$F : \mathcal{C}_1 \to \mathcal{C}_2$ between admissibility categories that intertwine
the respective reduction functors and preserve the admissibility predicates.

A functor $F$ between admissibility categories is \emph{admissibility-preserving} if
there exists a natural transformation
\[
  \phi : F \circ \mathfrak{R}_1 \Rightarrow \mathfrak{R}_2 \circ F
\]
making the reduction square commute, and if $F$ maps admissible morphisms in
$\mathcal{C}_1$ to admissible morphisms in $\mathcal{C}_2$.

The transport claims made in the main text correspond to the hypothesis that
admissibility-preserving functors exist between the admissibility categories determined
by each domain. Verifying this hypothesis in any specific case is a well-posed
mathematical problem and constitutes a research target within the Layer~V program.

% ================================================================
\section{Fiber Bundles, Projection, and Semantic Collapse}
\label{app:fiberbundle}
% ================================================================

This appendix develops the fiber-bundle interpretation of admissibility collapse,
giving a geometric formulation of the distinction between admissibility-preserving
compression and destructive projection.

\subsection{The projection setup}

Let $\pi : X \to M$ be a smooth surjection from a high-dimensional trajectory space
$X$ onto a lower-dimensional manifold $M$. Interpret $X$ as encoding the full
historical microstructure of a system, while $M$ encodes the compressed observable
structure accessible to an external observer or a downstream process. For each point
$m \in M$, the fiber $\pi^{-1}(m)$ consists of all microscopic states compatible with
macroscopic representation $m$.

A section $s : M \to X$ assigns to each macroscopic state a canonical microscopic
representative. The existence of globally consistent sections is the analog in this
geometric language of the existence of global sections in sheaf theory. When no global
section exists, locally consistent microstructure cannot be assembled into a globally
coherent representation.

\begin{definition}[Admissible projection]
  The projection $\pi : X \to M$ is \emph{admissible} if the following conditions hold.
  First, admissible reductions in $X$ project to continuous paths in $M$: if
  $\gamma : [0,1] \to X$ is an admissible trajectory then $\pi \circ \gamma$ is a
  continuous path in $M$. Second, admissible reductions contract fibers monotonically:
  if $x_0, x_1 \in X$ are connected by an admissible trajectory and $\pi(x_0) =
  \pi(x_1) = m$, then the fiber $\pi^{-1}(m)$ does not expand under the reduction.
  Third, invariant structure is transportable across local trivializations: for every
  local trivialization $\phi_i : \pi^{-1}(U_i) \to U_i \times F$ of the bundle, the
  invariant functor $\iota$ commutes with the transition maps $\phi_j \circ \phi_i^{-1}$
  on overlapping charts $U_i \cap U_j$.
\end{definition}

The third condition is precisely the cocycle condition for a fiber bundle with structure
group determined by the admissibility predicates. Its failure corresponds to the gluing
obstruction: locally admissible structure that cannot be assembled coherently.

\subsection{Curvature of the admissibility bundle}

When the admissibility bundle $\pi : X \to M$ admits a connection $\nabla$, its
curvature two-form $\Omega \in \Omega^2(M, \mathrm{End}(F))$ measures the failure of
parallel transport to be path-independent. Two paths $\gamma_1, \gamma_2$ in $M$ from
$m_0$ to $m_1$ produce holonomy operators $P_{\gamma_1}, P_{\gamma_2}$ on the fiber
$\pi^{-1}(m_0)$; their difference is encoded by the holonomy group and controlled by
$\Omega$.

In the admissibility framework, non-vanishing curvature on admissible paths corresponds
to path-dependence in admissibility transport: two reduction sequences beginning at the
same admissibility region may reach structurally inequivalent successors even if both
are individually admissible. Confluence is therefore the flatness condition on the
admissibility connection restricted to admissible paths: the curvature form $\Omega$
vanishes on pairs of tangent vectors determined by admissible reductions.

\subsection{Three failure modes as geometric pathologies}

Three distinct failure modes become precise in this geometric language.

Projection collapse occurs when distinct states $x_0, x_1 \in X$ satisfy
$\pi(x_0) = \pi(x_1)$ even though $x_0 \neq x_1$. Their admissibility distinction
is erased by the projection, and subsequent evolution over $M$ cannot distinguish
consequences that were distinguishable in $X$. This is the geometric form of lossy
compression: the map discards information that matters for future admissibility.

Fiber fragmentation occurs when the fiber $\pi^{-1}(m)$ splits into disconnected
components under reduction, so that globally consistent microscopic representatives
become unavailable even though local sections still exist. The system loses track
of the causal relationships between its own microstructures.

Topological hallucination occurs when local sections $s_i : U_i \to \pi^{-1}(U_i)$
exist for every open set in an open cover of $M$ and agree pairwise on overlaps
$U_i \cap U_j$, but no global section $s : M \to X$ exists. The first \v{C}ech
cohomology $\check{H}^1(M, \mathcal{F})$ is non-zero, where $\mathcal{F}$ is the
sheaf of local sections. This is the geometric characterization of hallucination as
a topological obstruction rather than as a local consistency failure.

\subsection{The map-territory instability as bundle degeneration}

When a system begins optimizing over $M$ while treating $M$ as identical to $X$, the
bundle structure degenerates. The admissibility predicates defined relative to $X$
become redefined relative to $M$, the connection on the bundle is replaced by a
degenerate flat connection, and the transition maps $\phi_j \circ \phi_i^{-1}$ become
trivial, collapsing the non-trivial fiber geometry into a product bundle and erasing
the admissibility-bearing structure of the fibers.

This is the geometric form of Goodhart's Law: optimizing a measure destroys its
correlation with the underlying structure it was intended to represent. The
admissibility framework implies that any representation subsequently subjected to
optimization pressure must preserve non-trivial fiber geometry from the outset.
Representations that are flat by construction cannot support coherent admissibility
transport under optimization.

% ================================================================
\section{Sheaf Cohomology, Semantic Coherence, and Hallucination}
\label{app:sheaf}
% ================================================================

This appendix develops the sheaf-cohomology interpretation of semantic coherence and
hallucination in detail, extending the Proposition in Section~\ref{sec:unification}
to a full cohomological analysis.

\subsection{Presheaves and sheaves on semantic manifolds}

Let $X$ be a topological space representing a semantic manifold and let
$\mathcal{U} = \{U_i\}_{i \in I}$ be an open cover of $X$. Let $\mathcal{F}$ be a
presheaf on $X$ with values in the category of semantic structures, so that
$\mathcal{F}(U)$ consists of locally consistent semantic assignments over $U$ and
the restriction maps $\rho_{UV} : \mathcal{F}(U) \to \mathcal{F}(V)$ for $V \subseteq U$
encode semantic specialization.

The presheaf $\mathcal{F}$ is a sheaf if two conditions hold. The locality condition
asserts that sections agreeing on every element of a cover are equal: if $s, t \in
\mathcal{F}(U)$ with $\rho_{U, U_i}(s) = \rho_{U, U_i}(t)$ for all $i$, then $s = t$.
The gluing condition asserts that a compatible family $\{s_i \in \mathcal{F}(U_i)\}$
satisfying $\rho_{U_i, U_i \cap U_j}(s_i) = \rho_{U_j, U_i \cap U_j}(s_j)$ for all
$i, j$ assembles into a unique global section $s \in \mathcal{F}(X)$ with
$\rho_{X, U_i}(s) = s_i$.

\subsection{\v{C}ech cohomology and obstruction theory}

For an open cover $\mathcal{U} = \{U_i\}_{i \in I}$, the \v{C}ech cochain groups are
\[
  \check{C}^n(\mathcal{U}, \mathcal{F})
  = \prod_{i_0 < i_1 < \cdots < i_n} \mathcal{F}(U_{i_0} \cap U_{i_1} \cap \cdots
  \cap U_{i_n}),
\]
with coboundary maps $\delta^n : \check{C}^n \to \check{C}^{n+1}$ defined by
alternating restriction maps in the standard way. The cohomology groups are
\[
  \check{H}^n(\mathcal{U}, \mathcal{F}) = \ker \delta^n / \operatorname{im} \delta^{n-1}.
\]
The zeroth group $\check{H}^0(X, \mathcal{F})$ is the group of global sections.
The first group $\check{H}^1(X, \mathcal{F})$ measures the obstruction to assembling
locally coherent patches into a global section: a one-cocycle $\{s_{ij} \in
\mathcal{F}(U_i \cap U_j)\}$ satisfying
\[
  s_{jk} - s_{ik} + s_{ij} = 0 \in \mathcal{F}(U_i \cap U_j \cap U_k)
\]
is a coboundary if and only if there exist $\{t_i \in \mathcal{F}(U_i)\}$ with
$s_{ij} = t_j - t_i$. Non-coboundary cocycles are the obstructions.

\begin{definition}[Semantic hallucination]
  A \emph{semantic hallucination} with respect to the cover $\mathcal{U}$ and
  presheaf $\mathcal{F}$ is a non-zero element of $\check{H}^1(X, \mathcal{F})$:
  an isomorphism class of \v{C}ech one-cocycles that cannot be written as coboundaries.
  Equivalently, it is a collection of locally compatible semantic patches that cannot
  be assembled into a globally coherent semantic structure.
\end{definition}

The characteristic phenomenology of hallucination follows directly. Each local patch
$s_i \in \mathcal{F}(U_i)$ is internally consistent. The pairwise compatibility
conditions are satisfied on every overlap. But the global section fails to exist
because the transition data form a non-trivial cohomology class. Local inspection
cannot detect the obstruction; only a global coherence check across triple overlaps
can reveal it.

\subsection{Admissibility as exactness and higher obstruction}

The vanishing of $\check{H}^1(X, \mathcal{F})$ is equivalent to exactness of the
sequence
\[
  0 \to \check{H}^0(X,\mathcal{F}) \to \prod_i \mathcal{F}(U_i)
  \xrightarrow{\delta^0} \prod_{i < j} \mathcal{F}(U_i \cap U_j).
\]
This exactness condition is the sheaf-theoretic formulation of semantic admissibility:
every locally admissible structure passing the compatibility test at overlaps must be
globally admissible. The admissibility framework asserts that a semantic system is
admissibility-sound if and only if $\check{H}^1(X,\mathcal{F}) = 0$.

The higher groups $\check{H}^n$ for $n \geq 2$ measure cascaded obstructions.
A non-trivial class in $\check{H}^2$ indicates failures of compatibility among triple
overlaps assembled from locally consistent pairwise data, corresponding in language
models to systems producing individually plausible facts that are pairwise consistent
but collectively contradictory through their logical implications.

% ================================================================
\section{Generalized Entropy-Constraint Dynamics}
\label{app:rsvp}
% ================================================================

This appendix develops a partial mathematical structure for RSVP-style entropy-constraint
dynamics, making the admissibility functional and variational equations as precise as
the current state of the framework permits.

\subsection{The RSVP field system}

Let $\Omega \subseteq \mathbb{R}^n$ be a compact region with smooth boundary, representing
the spatial domain of the plenum. The state at time $t$ is described by the scalar
coherence potential $\Phi : \Omega \times [0,T] \to \mathbb{R}$, the vector flow field
$\mathbf{v} : \Omega \times [0,T] \to \mathbb{R}^n$, and the entropy density
$S : \Omega \times [0,T] \to \mathbb{R}_{\geq 0}$.

The \emph{admissibility functional} is defined as
\[
  \mathcal{A}[\Phi, \mathbf{v}, S]
  = \int_\Omega \left(
    \alpha \, |\nabla \Phi|^2
    + \beta \, |\nabla \times \mathbf{v}|^2
    + \gamma \, S
    + \mu \, (\nabla \cdot \mathbf{v})^2
  \right) dV,
\]
where $\alpha, \beta, \gamma, \mu > 0$ are coupling constants. The four terms penalize
respectively: scalar gradient concentration, vortical flow structure, entropy density,
and compressional modes in $\mathbf{v}$.

Admissible trajectories satisfy two simultaneous constraints. The variational inequality
$\frac{d}{dt}\mathcal{A}[\Phi(t), \mathbf{v}(t), S(t)] \leq 0$ asserts that the
admissibility functional decreases along the trajectory. The local entropy production
constraint
\[
  \dot{S}(x,t) = \nabla \cdot \mathbf{J}_S(x,t) + \sigma(x,t) \geq 0,
\]
combined with the global stability condition
\[
  \int_\Omega \sigma(x,t) \, dV < \dot{\Sigma}_{\mathrm{crit}},
\]
prevents supercritical entropy production.

\subsection{Euler--Lagrange equations for admissible attractors}

Setting the first variation of $\mathcal{A}$ to zero yields the Euler--Lagrange system.
Varying with respect to $\Phi$ gives
\[
  -2\alpha \, \Delta \Phi = \lambda_\Phi,
\]
where $\lambda_\Phi$ is a Lagrange multiplier enforcing the entropy coupling. Varying
with respect to $\mathbf{v}$ gives, for divergence-free flow fields,
\[
  2\beta \, \Delta \mathbf{v} = \boldsymbol{\lambda}_v.
\]
An \emph{admissible attractor} $(\Phi^\ast, \mathbf{v}^\ast, S^\ast)$ satisfies these
equations with $\lambda_\Phi = \gamma$ and $\boldsymbol{\lambda}_v = \mathbf{0}$:
\[
  \Delta \Phi^\ast = -\frac{\gamma}{2\alpha}, \qquad
  \Delta \mathbf{v}^\ast = \mathbf{0},
\]
together with $\int_\Omega \sigma^\ast \, dV < \dot{\Sigma}_{\mathrm{crit}}$.

\subsection{Stability analysis and the critical coupling}

To analyze stability of $(\Phi^\ast, \mathbf{v}^\ast, S^\ast)$ under perturbations
$(\delta\Phi, \delta\mathbf{v}, \delta S)$, expand $\mathcal{A}$ to second order:
\[
  \delta^2 \mathcal{A}
  = \int_\Omega \left(
    \alpha \, |\nabla(\delta\Phi)|^2
    + \beta \, |\nabla \times (\delta\mathbf{v})|^2
    + \gamma \, \delta S
  \right) dV.
\]
Since $\alpha, \beta > 0$, the first two terms are positive-definite. The stability
condition reduces to $\gamma \int_\Omega \delta S \, dV > 0$ for admissible entropy
perturbations. Since admissibility requires $\delta S \geq 0$ pointwise, the condition
is equivalent to $\gamma > 0$, which holds by assumption. Every admissible attractor
is therefore a local minimum of $\mathcal{A}$ within the admissible sector.

The parameter $\lambda_c$ of the main text can be identified with
\[
  \lambda_c = \frac{\gamma}{2\alpha},
\]
which controls the magnitude of the uniform Laplacian of $\Phi^\ast$ and sets the
characteristic length scale of the coherence field. For $\lambda_c \approx 0.42$,
the coherence field has a characteristic wavelength consistent with stable attractor
formation in the model. This is the sense in which $\lambda_c$ is a heuristic scaling
parameter: it encodes the geometry of admissible attractors within the current
variational framework, not an empirical universal constant.

% ================================================================
\section{Homotopy-Theoretic Reading of Identity Persistence}
\label{app:homotopy}
% ================================================================

This appendix develops the homotopy-theoretic interpretation of identity under
transformation, arguing that the admissibility framework requires replacing extensional
identity with homotopical identity.

\subsection{The inadequacy of extensional identity}

Classical symbolic systems identify objects by their extensional properties: two objects
are the same if they have the same parts or satisfy the same predicates. This
identification is static and ahistorical. The admissibility framework requires a
different notion because condition~(iii) of Definition~\ref{def:R} asserts that history
is strictly extended by each reduction: the same extensional state reached by two
different admissibility trajectories may carry different structural consequences.

Formally, let $x_0, x_1 \in X$ be two points in the trajectory space with
$\pi(x_0) = \pi(x_1) = m \in M$, so that they are extensionally identical at the
level of the compressed representation. The admissibility framework distinguishes them
if the homotopy classes of their approach trajectories differ: the path taken to reach
$x_0$ and the path taken to reach $x_1$ may have deposited different admissibility
history, making the two points structurally inequivalent despite extensional coincidence.

\begin{definition}[Admissible identity path]
  An \emph{admissible identity path} from $x_0$ to $x_1$ is a continuous map
  $\gamma : [0,1] \to X$ with $\gamma(0) = x_0$ and $\gamma(1) = x_1$ such that every
  intermediate state $\gamma(t)$ is admissible for all $t \in [0,1]$.
\end{definition}

Two states are \emph{admissibly homotopic} if there exists an admissible identity path
between them. The admissibly homotopic equivalence relation partitions $X$ into
connected components of the admissible locus $X_{\mathsf{adm}}$.

\subsection{The fundamental group of the admissible locus}

The fundamental group $\pi_1(X_{\mathsf{adm}}, x_0)$ encodes homotopy classes of
admissible loops based at $x_0$. In the Spherepop interpretation, a non-trivial loop
corresponds to a sequence of membrane collapses that returns the bubble system to its
original topology but not its original state: the history has changed even though the
extensional description has not.

In the cognitive interpretation, a non-trivial loop in the admissible conceptual space
corresponds to a thought trajectory returning to an apparently familiar concept while
arriving from a different direction, encoding a different set of analogical associations.
Hofstadter's strange loops are precisely the non-trivial elements of
$\pi_1(X_{\mathsf{adm}}, x_0)$.

\subsection{Higher homotopy groups and coherence conditions}

The higher homotopy groups $\pi_n(X_{\mathsf{adm}}, x_0)$ for $n \geq 2$ encode
higher-dimensional coherence conditions. Vanishing of $\pi_2$ asserts that any two
admissible homotopies between the same pair of paths can themselves be admissibly
homotoped to each other, a condition analogous to Church--Rosser: just as Church--Rosser
asserts that reduction sequences from a common source can be joined at a common target,
vanishing of $\pi_2$ asserts that homotopies between common paths can be joined by a
coherent higher homotopy.

The connection between higher homotopy groups and higher cohomology is made precise
by the Hurewicz theorem
\[
  h : \pi_n(X_{\mathsf{adm}}, x_0) \to H_n(X_{\mathsf{adm}}; \mathbb{Z}),
\]
which is an isomorphism when $X_{\mathsf{adm}}$ is $(n-1)$-connected. A complete
admissibility theory would compute these invariants for each domain and compare them,
establishing whether the admissibility structures in different domains are
homotopy-equivalent or related only at the level of the fundamental group.

The Homotopy Type Theory program provides a natural foundational setting for this
program. In HoTT, the identity type $\mathrm{Id}_A(x,y)$ is not a proposition but a
type whose inhabitants are identity witnesses, that is, paths from $x$ to $y$. The
admissible identity path of the definition above is precisely such a term.
Admissibility collapse corresponds to concatenation of identity witnesses, and the
history $H$ of a bubble in Spherepop is the accumulated data of all identity witnesses
traversed. A future formalization of Spherepop may therefore find its natural home in
cohesive homotopy type theory, which adds a geometric modality capable of expressing
the spatial structure of bubble membranes.

% ================================================================
\section{Non-Vacuity Criterion and Boundaries of the Framework}
\label{app:nonvacuity}
% ================================================================

This appendix develops the non-vacuity criterion for the admissibility framework in
detail, addressing the objection that any constrained process might instantiate the
generalized reduction operator.

\subsection{The vacuity problem}

The vacuity problem for any sufficiently general framework is that the framework's
conditions may be satisfiable by every process in the intended domain, rendering the
framework empirically empty. For the admissibility principle, the risk is real:
Definition~\ref{def:R} specifies four conditions but these are individually weak enough
that many processes satisfy them trivially. A thermostat satisfies all four under a
liberal reading, yet it is clearly not the intended kind of system.

The resolution is to augment Definition~\ref{def:R} with conditions that are jointly
strong enough to exclude trivial instantiations.

\subsection{Strong admissibility: the augmented criterion}

\begin{definition}[Strong admissibility]
  A process $\mathcal{P}$ satisfies \emph{strong admissibility} if it satisfies all
  four conditions of Definition~\ref{def:R} and the following three augmentations hold.
  First, the admissibility predicate $\mathsf{Adm}$ recursively restructures the future
  trajectory space: the set of admissible perturbations $\{\Delta : \mathsf{Adm}(\Bcal_i,
  \Delta)\}$ changes in a history-dependent way not determined by $\Bcal_i$ alone but
  by the full history $H(\Bcal_i)$.
  Second, the invariant $\iota(\Bcal_i)$ is transportable, meaning that downstream
  processes can use $\iota(\Bcal_{i+1})$ as a constraint on their own admissibility
  without accessing the full history $H(\Bcal_{i+1})$.
  Third, the system is path-sensitive: there exist pairs of histories $H$ and $H'$
  leading to the same extensional state $\Bcal$ but determining different future
  admissibility sets $\{\Delta : \mathsf{Adm}((\Bcal, H), \Delta)\} \neq \{\Delta :
  \mathsf{Adm}((\Bcal, H'), \Delta)\}$.
\end{definition}

These three augmentations together constitute a significant strengthening. The first
condition excludes all Markovian processes: in a Markov process the future transition
kernel depends only on the current state, so the admissibility predicate cannot be
recursively restructured by history. The second condition excludes processes that
propagate invariants only internally without making them available for downstream use.
The third condition excludes all path-independent processes, including thermostats,
lookup tables, and arbitrary dissipative systems.

\subsection{Canonical non-instances and their analysis}

A thermostat fails the first condition because its future behavior is determined
entirely by the current temperature reading, not the sequence of prior readings.
The admissibility predicate does not change as a function of the history of prior
adjustments. A thermostat instantiates weak admissibility at best, an output filter
rather than a recursive state-space reorganizer.

A lookup table fails all three augmented conditions. Its future outputs are determined
by current inputs alone, the table values are invariant under any sequence of queries,
and the history of prior queries has no effect on future outputs.

A memoryless Markov chain fails the first and third conditions: future transition
probabilities are determined by the current state alone, and no history-dependent
restructuring of the state space occurs.

A physically dissipative system such as a glass shattering fails the second condition:
no transportable invariant survives the transformation in a form usable by downstream
processes. The shards do not carry the organizational invariant of the intact glass
forward in any operationally meaningful way.

\subsection{Biological development as a positive exemplar}

The development of a multicellular organism from a fertilized egg is the canonical
positive instance of strong admissibility; it is treated at length in
Section~\ref{sec:lineage} through the computational biology of multimodal
lineage tracing (Wang, He, \& Hu 2026). Here the essential formal features
are summarized. At each stage, the admissibility predicate
governing which cell fates are accessible is determined not merely by the cell's current
gene expression state but by its full developmental history: the lineage of prior
differentiation events, the morphogen gradients experienced, and the positional
information accumulated through tissue interactions. This is recursive state-space
restructuring in precisely the sense required by the first augmented condition.

The transportable invariant is positional and lineage identity: the information that a
cell's descendants will use to determine their own admissible developmental options.
This invariant is communicated through signaling molecules, transcription factor profiles,
and epigenetic marks, making it externally accessible to neighboring cells without
requiring access to the full developmental history. This satisfies the second augmented
condition.

Finally, the system is strongly path-sensitive: two cells with identical current gene
expression but different developmental histories will have different admissible futures
even if their present transcriptional states are extensionally indistinguishable.
This is precisely the third augmented condition. Biological development is therefore
both a positive exemplar of strong admissibility and an existence proof that naturally
occurring systems can satisfy all augmented conditions simultaneously.

% ================================================================
\section{The Three-Level Coarse-Graining: Syntax, Topology,
  Thermodynamics}
\label{app:threelevel}
% ================================================================

\subsection{The coarse-graining hierarchy}

A persistent objection to the admissibility framework is that it appears to dissolve
the distinctions between fundamentally different ontological categories. Symbolic
rewriting systems, topological deformations, and thermodynamic evolutions are not
members of the same mathematical universe, and identifying them risks category errors
of a severe kind.

The resolution is that the admissibility framework does not identify these levels but
asserts that the same admissibility grammar reappears at each level of coarse-graining.
The progression is a hierarchy:
\[
  \text{syntax} \;\xrightarrow{\text{geometric realization}}\; \text{topology}
  \;\xrightarrow{\text{thermodynamic limit}}\; \text{thermodynamics},
\]
in which each arrow represents a coarse-graining operation preserving admissibility
structure while suppressing detail.

At the syntactic level, admissibility is determined by typing rules, capture-avoidance
conditions, and the structure of explicit substitutions. The admissibility predicate
is decidable: type-checking is computable, and the invariant is the typing judgement.

At the topological level, syntactic objects are replaced by geometric ones: terms
become bubbles, substitution becomes membrane collapse, and the reduction relation
becomes a continuous deformation of topological spaces. Admissibility is no longer
decidable in general but it is geometrically natural. The invariant is topological:
the homotopy type of the admissible locus, its fundamental group, the higher homotopy
groups, and the sheaf cohomology.

At the thermodynamic level, topological objects are replaced by field configurations:
bubbles become regions of the entropy-potential field, membrane collapse becomes
gradient-constrained field relaxation, and reduction becomes a variational flow.
Admissibility is a condition on the entropy production rate and the admissibility
functional. The invariant is thermodynamic: the free energy, the entropy production
history, and the admissible attractor basin.

\subsection{The mesoscale necessity theorem}

A structural consequence of this three-level hierarchy is the necessity of mesoscale
organization. At the syntactic level, admissibility predicates operate on terms of
finite size. At the thermodynamic level, admissibility conditions are expressed as
integrals over extended regions. Neither level alone can support coherent admissibility
transport across the full range of scales relevant to cognitive, physical, and semantic
systems.

The topological level is the necessary intermediary. It provides geometric objects of
finite extent but with non-trivial topological structure—bubble complexes, homotopy
cells, sheaf stalks—that can simultaneously encode syntactic structure through their
combinatorial topology and thermodynamic structure through their volume and boundary
geometry.

The admissibility principle therefore predicts that stable computation, cognition, and
physical organization require persistent mesoscale constraint structures. In turbulence,
Reynolds stresses encode mesoscale admissibility information that is neither purely
microscopic (individual fluid particle trajectories) nor purely macroscopic (mean flow
statistics). In cognition, conceptual chunks are mesoscale attractors between neuronal
microstates and explicit symbolic reasoning. In language models, embedding
representations occupy the mesoscale between token-level statistics and global semantic
coherence. Any theory that eliminates mesoscale organization should fail to produce
coherent admissibility transport: this is a nontrivial and potentially falsifiable
prediction of the framework.

\subsection{The thermodynamic cost of syntactic irreversibility}

The relationship between the syntactic and thermodynamic levels can be made partially
precise through computational thermodynamics. Landauer's principle asserts that any
logically irreversible computation must dissipate at least $k_B T \ln 2$ of energy
per bit erased, where $k_B$ is the Boltzmann constant and $T$ is the temperature of
the thermal environment.

In the admissibility framework, condition~(iii) of Definition~\ref{def:R} asserts that
each reduction strictly extends the history. This means that the admissibility framework
is inherently logically irreversible in the Landauer sense: each collapse erases the
pre-collapse state from the admissible futures, and this erasure has a minimum
thermodynamic cost proportional to the information content of the erased pre-collapse
state.

The total thermodynamic cost of a reduction sequence $\Bcal_0 \to \Bcal_1 \to \cdots
\to \Bcal_n$ is at least
\[
  W \geq n \cdot k_B T \ln 2 + \sum_{k=0}^{n-1} k_B T \cdot \log_2 |\Bcal_k|,
\]
where $|\Bcal_k|$ is the information content of the pre-collapse state at step $k$.
The RSVP entropy production limit $\dot{\Sigma}_{\mathrm{crit}}$ can therefore be
interpreted as an upper bound on the rate at which Landauer-irreversible operations
can be performed: a system exceeding this limit produces entropy faster than its
constraint geometry can absorb, leading to structural destabilization.

This connection identifies a specific and tractable research target for Layer~V:
deriving $\dot{\Sigma}_{\mathrm{crit}}$ from Landauer's principle applied to
admissibility-irreversible reductions. The derivation would require a precise
information-theoretic account of the information content of admissibility regions and
the Landauer cost of each collapse, which in turn requires the formal Spherepop
specification identified in Section~\ref{sec:spherepop} as the primary open problem.
The three formal gaps — Spherepop axiomatization, categorical embedding of $\mathcal{R}$,
and thermodynamic derivation of $\dot{\Sigma}_{\mathrm{crit}}$ — are therefore not
independent: progress on any one of them directly advances the others.

\end{document}
